\section{The Extended Hecke Action}\label{sec:extendedheckestacks}
Continue to fix $G$ a quasi-split connected reductive group over $F$, with a fixed finite central subgroup $Z$.
\subsection{Hecke stacks}\label{sec:hecke-stacks}
Let $S = \Spa(R,R^+) \in \Perfd$. Given a map $x : S \to \Div^1$ corresponding to an untilt $\Spa(R^{\#}, R^{\#,+})$, we denote by 
$$\D_x = \Spec(\Bdr^{+}(R^{\sharp}))$$
and by 
$$\D_x^{\circ} = \Spec(\Bdr(R^{\sharp})).$$
More generally for $\underline{x} : S \to (\Div^1)^d$, we denote by $\D_{\underline{x}} \xrightarrow{\iota_{\underline{x}}} X_{S}$ the spectrum of the completion of the corresponding (degree $d$) closed Cartier divisor on $X_S$ and by $\D_{\underline{x}}^{\circ} = \D_{\underline{x}} - \underline{x} \xrightarrow{\iota_{\underline{x}}^{\circ}} X_{S}$, see \cite[Section VI.1]{Geometrization}.

We first recall the local and global Hecke stacks used in \cite{Geometrization}.

\begin{defi}[Local Hecke stack with $I$ legs]
Fix a finite set $I$ equipped with an ordered partition $I = \bigsqcup_{i=0}^{k} I_{i}$.
Define the functor $\tn{Hck}_{G}^{\tn{loc},(I_{0}, \dots, I_{k})}$ on $\Perfd$ by sending $S$ to the following data:
\begin{enumerate}
\item{A map $S \xrightarrow{\underline{x} = (\underline{x}_{i})_{i=0}^{k}}(\Div^{1})^{I}$;}
\item{A $k+1$-tuple $(\cP_{i})_{i=0}^{k}$ of $G$-torsors on $\bbD_{\underline{x}}$;}
\item{A sequence $(f_{i})_{i=0}^{k-1}$ of isomorphisms of $G$-torsors
\begin{equation*}
    \cP_{i}|_{\bbD_{\underline{x}} \setminus \underline{x}_{i}} \xrightarrow{\sim} \cP_{i+1}|_{\bbD_{\underline{x}} \setminus \underline{x}_{i}}.
\end{equation*}
}
\end{enumerate}
\end{defi}

\begin{defi}[Global Hecke stack with $I$ legs]
Fix a finite set $I$ equipped with an ordered partition $I = \bigsqcup_{i=0}^{k} I_{i}$.
Define the functor $\tn{Hck}_{G}^{(I_{0}, \dots, I_{k})}$ on $\Perfd$ by sending $S$ to the following data:
\begin{enumerate}
\item{A map $S \xrightarrow{\underline{x} = (\underline{x}_{i})_{i=0}^{k}}(\Div^{1})^{I}$;}
\item{A $k+1$-tuple $(\cP_{i})_{i=0}^{k}$ of $G$-torsors on $X_{S}$;}
\item{A sequence $(f_{i})_{i=0}^{k-1}$ of isomorphisms of $G$-torsors
\begin{equation*}
    \cP_{i}|_{X_{S} \setminus \Gamma_{\underline{x}_{i}}} \xrightarrow{\sim} \cP_{i+1}|_{X_{S} \setminus \Gamma_{\underline{x}_{i}}}.
\end{equation*}
}
\end{enumerate}
\end{defi}

We now give the first ``extended'' variant of these objects. It will turn out that, although this is the most natural definition of an ``extended local Hecke stack,'' it is not quite enough for our full purposes.

The analogue of Lemma \ref{lem:inertialhom} does not hold for $\bbD_{\underline{x}}$, so we need to make the following definition:


\begin{defi}\label{defi:bandloc}
    For $S = \Spa(R,R^{+})$ in $\Perfd$ and $\underline{x} \in (\Div^{1})^{I}(S)$, we say that $\cP$ a $G_{S}$-torsor on $[\iota_{\underline{x}}^{*}\mc{T}_{S}^{\loc}/\tilde{t}_{S}]$ is \textbf{band-localized} if $\lambda_{\cP} \in \Hom_{\bbD_{\underline{x}}}(u_{S}, Z_{G,S})$ extends (uniquely) to a morphism in $\Hom_{X_{S}}(u_{S}, Z_{G,S})$. Note that $\lambda_{\cP}$ factors through a fixed finite subgroup $Z \subseteq Z_{G}$ if and only if this global extension does.
\end{defi}

\begin{defi}\label{defi:hcknaive}
    Fix a finite set $I$ equipped with an ordered partition $I = \bigsqcup_{i=0}^{k} I_{i}$.
Define the functor $\tn{Hck}_{G}^{e,(I_{i}),\tn{naive}}$, resp. $\tn{Hck}_{G}^{e,\tn{loc},(I_{i}),\tn{naive}}$ on $\Perfd$ by sending $S$ to the following data:
\begin{enumerate}
\item{A map $S \xrightarrow{\underline{x} = (\underline{x}_{i})_{i=0}^{k}}(\Div^{1})^{I}$;}
\item{A $k+1$-tuple $(\cP_{i})_{i=0}^{k}$ of $G$-torsors on $[\mc{T}_{S}/\tilde{t}_{S}]$, resp. band-localized $G$-torsors on $[\iota_{\underline{x}}^{*}\mc{T}_{S}/\tilde{t}_{S}]$;}
\item{A sequence $(f_{i})_{i=0}^{k-1}$ of isomorphisms of $G$-torsors
\begin{equation*}
    \cP_{i}|_{[\mc{T}_{S}/\tilde{t}_{S}] \times_{X_{S}} X_{S} \setminus \Gamma_{\underline{x}_{i}}} \xrightarrow{\sim} \cP_{i+1}|_{[\mc{T}_{S}/\tilde{t}_{S}] \times_{X_{S}} X_{S} \setminus \Gamma_{\underline{x}_{i}}},
\end{equation*}
resp. an isomorphism of $G$-torsors
\begin{equation*}
    \cP_{i}|_{[\iota_{\underline{x}}^{*}\mc{T}_{S}/\tilde{t}_{S}] \times_{\bbD_{\underline{x}}} \bbD_{\underline{x}} \setminus \underline{x}_{i}} \xrightarrow{\sim} \cP_{i+1}|_{[\iota_{\underline{x}}^{*}\mc{T}_{S}/\tilde{t}_{S}] \times_{\bbD_{\underline{x}}} \bbD_{\underline{x}} \setminus \underline{x}_{i}}.
\end{equation*}
}
\end{enumerate}
One can also define $\tn{Hck}_{Z \subset G}^{e,(I_{i}),\tn{naive}}$ and $\tn{Hck}_{Z \subset G}^{e,\tn{loc},(I_{i}),\tn{naive}}$ by taking the above definitions but insisting that all $\lambda_{\cP}$ factor through the fixed subgroup $Z$.
\end{defi}
All of the above four stacks are equipped with natural convolution operations as defined in \cite[\S VI.8]{Geometrization}. 

 \begin{rque}
     We note that, for any $S = \Spa(R,R^{+}) \in \Perfd$ and $\underline{x} \in (\tn{Div}^{1})^{I}(S)$, for the local part of Definition \ref{defi:hcknaive} to make sense we are implicitly viewing the $\tilde{t}_{S}$-torsor $\mc{T}_{S}^{\loc}$, which we defined as the pullback to the formal disk around $\underline{x}$ of $\mc{T}_{S}$, as an \' {e}tale $\tilde{t}_{S}$-torsor over $\Spec(B^{+}_{(\Div^{1})^{I}}(R^{\sharp}))$ (see \cite[\S VI.1]{Geometrization} for the precise definition of this ring), even though $\mc{T}_{S}^{\loc}$ is a priori an \'{e}tale $\tilde{t}_{S}$ torsor on the associated adic space. 
     
     Since $\mc{T}_{S}^{\loc}$ is pulled back from a (unique) torsor on $\Spec(B^{+}_{(\Div^{1})^{I}}(R^{\sharp}))$---this follows from the explicit description of $\mc{T}_{S}$ as coming from finite Galois extensions of $F$ and isocrystals---the above maneuver is harmless, and we use it throughout this paper. For completeness, the claimed algebraic torsor is the one determined by the norm-compatible system (across all finite Galois $E/F$) of line bundles given by taking the sections of the line bundle $\mc{E}(\mathbb{L}_{E})$ on $X^{(E)}_{S}$ at the completed neighborhood of the graph of $\underline{x}$, viewed as a module over the ring $B^{+}_{(\Div^{1})^{I}}(R^{\sharp})$.
 \end{rque}


\begin{defi}\label{defi:hom} Denote by $\underline{\Hom}_{X}(u,Z)$ the sheaf on $\Perfd$ sending $S$ to $\Hom_{X_{S}}(u_{S},Z_{S})$. 
\end{defi}
We deduce that there are locally constant (by the proof of Lemma \ref{lem:inertialhom}) morphisms
\begin{equation*}
    \Hck_{G}^{e,(I_{i}),\tn{naive}}, \Hck_{G}^{e,\loc,(I_{i}),\tn{naive}}\to \underline{\Hom}_{X}(u,Z),
\end{equation*}
and define, for a fixed $\lambda \in \Hom_{F}(u,Z_{G})$, the substacks $\Hck_{G}^{e,(I_{i}),\tn{naive},\lambda}, \Hck_{G}^{e,\loc,(I_{i}),\tn{naive},\lambda}$ to be the preimages of $\lambda$ under these morphisms (one could rephrase this condition to match the statement of Proposition \ref{prop:BunGestructure} by insisting that $\lambda_{\cP_{i,\bar{s}}} = \lambda$ for all geometric points $\bar{s}$ of $S$). 

\begin{lem}\label{lem:naiveHck}
    For a fixed $\lambda \in \Hom_{F}(u,Z_{G})$ and a fixed $b \in B_{e}(G)_{\tn{basic}}$ whose restriction to $u$ is $\lambda$, we can twist by $b$ to obtain isomorphisms
    \begin{equation*}
        \Hck_{G}^{e,(I_{i}),\tn{naive},\lambda} \xrightarrow{\sim} \Hck_{G_{b}}^{{(I_{i})}}
    \end{equation*}
    and 
    \begin{equation*}
        \Hck_{G}^{e,\loc,(I_{i}),\tn{naive},\lambda} \xrightarrow{\sim} \Hck_{G_{b}}^{\loc,{(I_{i})}}.
    \end{equation*}
\end{lem}

\begin{proof}
    We have already seen in Proposition \ref{prop:BunGestructure}.(3) that twisting gives an isomorphism $\Bun_{G_{b}} \to \Bun_{G}^{e,\lambda}$, and this twisting bijection holds on the local disks $\bbD_{\underline{x}}$ as well and is functorial on isomorphisms (it is given by applying the functor $- \times^{G_{b}} \cP_{b}$ for the fixed $G$-torsor $\cP_{b}$ over $[\mc{T}_{C}/\tilde{t}_{C}]$ corresponding to $b$, where $\cP_{b}$ is viewed as a $(G_{b},G)$-bitorsor on $[\mc{T}_{C}/\tilde{t}_{C}]$).
\end{proof}

We deduce from combining Lemma \ref{lem:naiveHck} with \cite[Section IX]{Geometrization} that the stacks $\Hck_{G}^{e,(I_{i}),\tn{naive}}$ and $\Hck_{G}^{e,\loc,(I_{i}),\tn{naive}}$ give, for any finite set $I$ and $V \in \tn{Rep}((\widehat{G} \rtimes Q)^{I})$ (where $Q$ is a finite quotient of $\Weil_{F}$ through which the action on $\widehat{G}$ factors), a Hecke operator
\begin{equation*}
    T^{e}_{V} \colon \mc{D}_{\tn{lis}}(\Bun_{G}^{e},\Lambda) \to \mc{D}_{\tn{lis}}(\Bun_{G}^{e},\Lambda)^{B\Weil_{F}^{I}},
\end{equation*}
as in \cite[Theorem IX.0.1]{Geometrization}, which satisfies all of the properties as in the operators $T_{V}$ \emph{loc. cit.} 

We can therefore deduce immediately from \cite[Corollary X.1.3]{Geometrization}:

\begin{prop}\label{prop:naiveSA}
    There is a natural compactly supported $\Lambda$-linear action of $\tn{Perf}(Z^{1}(\Weil_{F},\widehat{G})_{\Lambda}/\widehat{G})$ on $\mc{D}_{\tn{lis}}(\Bun_{G}^{e},\Lambda)^{\omega}$ which is compatible with the Hecke action and gives, functorially in finite sets $I$, exact $\tn{Rep}_{\Lambda}(Q^{I})$-linear functors
    \begin{equation*}
        \tn{Rep}_{\Lambda}((\widehat{G} \rtimes Q)^{I}) \to \End_{\Lambda}(\mc{D}_{\tn{lis}}(\Bun_{G}^{e},\Lambda)^{\omega})^{B\Weil_{F}^{I}},
    \end{equation*}
    where $Q$ denotes a finite quotient of $\Weil_{F}$ through which the action on $\widehat{G}$ factors and the superscript $\omega$ denotes $\mc{D}_{\tn{lis}}(\Bun_{G}^{e},\Lambda)^{\omega} \subseteq \mc{D}_{\tn{lis}}(\Bun_{G}^{e},\Lambda)$ the subcategory of compact objects.
\end{prop}

To summarize the situation, the above proposition and the resulting spectral action, are obtained from the theory of \cite{Geometrization} by simply decomposing
\begin{equation}\label{eq:naivedecomp}
 \Hck_{G}^{e,(I_{i}),\tn{naive}} = \bigsqcup_{\lambda \in \Hom_{F}(u,Z)}  \Hck_{G}^{e,(I_{i}),\tn{naive},\lambda} \xrightarrow{\sim}  \bigsqcup_{\lambda \in \Hom_{F}(u,Z)}  \Hck_{G_{b_{\lambda}}}^{(I_{i})}
\end{equation}
(and its local analogue) using Lemma \ref{lem:naiveHck}.

In order to construct the full ``extended spectral action'' from \cite[Conjecture 12.8]{Fargues22}, one needs a version of the Hecke stacks $\Hck_{G}^{e,(I_{i}),\tn{naive}}$, $\Hck_{G}^{e,\loc,(I_{i}),\tn{naive}}$ which allows ``intertwining'' between different $\Hom_{F}(u,Z)$--components of $\Bun_{G}^{e}$.

We conclude with a technical lemma which will be needed later:

\begin{lem}\label{lem:Grgerbe}
    The action of $L^{+}G$ on $\tn{Gr}_{G/Z,\Spd(C)}$ has the same orbits as that of $L^{+}(G/Z)$.
\end{lem}

\begin{proof}
    In mixed characteristic this follows from the fact that $L^{+}G \to L^{+}(G/Z)$ is surjective on geometric points $\Spa(C')$, since if $x \colon \Spd(C') \to \Spd(C) \xrightarrow{y} \Div^{1}$ is the point induced by $y$ then $H^{1}_{\tn{\'{e}t}}(\bbD_{x}, Z) = 0$ because $\Bdr^{+}(C^{'\sharp})$ is a strictly Henselian local ring with algebraically closed residue field.

    In the general case, observe that for any geometric point $\Spa(C')$ we have the commutative diagram
    \[ 
\begin{tikzcd}
    L^{+}T(C') \arrow{r} \arrow{d} & L^{+}(T/Z)(C') \arrow{r} \arrow{d} & H^{1}_{\tn{fppf}}(\bbD_{x}, Z) \arrow["\tn{id}"]{d} \\
     L^{+}G(C') \arrow{r} & L^{+}(G/Z)(C') \arrow{r} & H^{1}_{\tn{fppf}}(\bbD_{x}, Z),
\end{tikzcd}
    \]
    which implies that the morphism $L^{+}T \backslash L^{+}(T/Z) \to L^{+}G \backslash L^{+}(G/Z)$ is surjective. 

   Recall that, by the Cartan decomposition, any point $[g] \in \tn{Gr}_{G/Z}(C')$ lies in the same $L^{+}(G/Z)(C')$-orbit as some $t \in LT(C')$. We then obtain the result by writing any $g \in L^{+}(G/Z)$ as $g = g't'$ for $g' \in L^{+}G(C')$ and $t' \in L^{+}(T/Z)(C')$ and observing that
   \begin{equation*}
       [gt] = [g't't] = [g'tt'] = [g't],
   \end{equation*}
giving the equality of the orbits.
\end{proof}

\subsection{Extended Hecke stacks}\label{sec:extendedheckedef}
We now define local and global Hecke stacks adapted to the extended moduli of bundles $\Bun_{G}^{e}$. 

    Applying the functor $\Hom_{F}(-,\mathbf{T})$ to the sequence 
    \begin{equation*}
1 \to u \to \tilde{t} \to t \to 1
    \end{equation*}
    and taking a fibered product yields the short exact sequence
    \begin{equation}\label{eq:HomZSES}
        1 \to \Hom_{F}(t, \mathbf{T}) \to \Hom_{F}(\tilde{t}, \mathbf{T}) \times_{\Hom_{F}(u,\mathbf{T})} \Hom_{F}(u,Z) \to \Hom_{F}(u,Z) \to 1,
    \end{equation}
    where the surjectivity comes from applying Shapiro's Lemma and using that $\tn{Ext}^{1}_{\Z}(X^{*}(\mathbf{T}), X^{*}(t)) = 0$.
    We fix, for once and for all, a normalized (set-theoretic) section 
\begin{equation*}
    s \colon \Hom_{F}(u,Z_{G}) \to \Hom_{F}(\tilde{t}, \mathbf{T}) \times_{\Hom_{F}(u,\mathbf{T})} \Hom_{F}(u,Z_{G}),
\end{equation*}
determining a section 
\begin{equation*}
s \colon \Hom_{F}(u,Z) \to \Hom_{F}(\tilde{t}, \mathbf{T}) \times_{\Hom_{F}(u,\mathbf{T})} \Hom_{F}(u,Z)
\end{equation*}
of the right-hand map in \eqref{eq:HomZSES} with corresponding normalized $2$-cocycle 
\begin{equation*}
z \in Z^{2}(\Hom_{F}(u,Z),\Hom_{F}(t, \mathbf{T})).
\end{equation*}
By Lemma \ref{lem:inertialhom} we get a section $\Hom_{X_{S}}(u_{S},Z_{G,S}) \to \Hom_{X_{S}}(\tilde{t}_{S},\mathbf{T}_{S})$ for all $S \in \Perfd$ as well, which we also denote by $s$.
    

As a notational convention, we set $\bbD_{y} =: \bbD_{C}$, and for a fixed $S = \Spa(R,R^{+}) \to \Spa(C)$ in $\Perfd_{C}$ with corresponding morphism $\Spa(R^{\sharp},R^{\sharp,+}) \xrightarrow{x} \Spa(C^{\sharp},C^{\sharp,+})$, we set $\bbD_{S} = \bbD_{y \circ x}$ and $\bbD_{S}^{\circ} = \bbD_{y \circ x}^{\circ}$. We also set $X_{S}^{\circ}:= X_{S} \setminus \Gamma_{y \circ x}$ and (recall the $t_{S}$-torsor $\mc{T}_{S}$ over $X_{S}$ from Definition \ref{defi:Tdef}) $\mc{T}_{S}^{\tn{loc}}$, $\mc{T}_{S}^{\tn{loc},\circ}$, or $\mc{T}_{S}^{\circ}$  the pullback of $\mc{T}_{S}$ to $\bbD_{S}$, $\bbD_{S} ^{\circ}$, or $X_{S}^{\circ}$, respectively.

Note that, for any affinoid $\Spa(R,R^{+}) = S$ in $\Perfd_{C}$, there is a map 
\begin{equation}\label{eq:gerbesm}
[\mc{T}^{\tn{loc}}_{S}/\tilde{t}_{S}] \to [\bbD_{S}/\tilde{t}_{S}]
\end{equation}
induced by the structure map $\mc{T}_{S}^{\tn{loc}} \to \bbD_{S}$, similarly for $[\mc{T}_{S}/\tilde{t}_{S}] \to [X_{S}/\tilde{t}_{S}]$.

\begin{lem}\label{lem:equivtors} Let $X$ be a site, $A$ an abelian sheaf of groups over $X$, and $H$ a sheaf of groups over $X$. Defining an $H$-torsor on $[X/A]$ is equivalent to defining an $A$-equivariant $H$-torsor on $X$. 
\end{lem}

\begin{proof}
    See \cite[Remark 2.6]{Shin21}.
    \end{proof}

    For an $H$-torsor $\mc{U}$ on a site $X$ and $\Q$ a $Z \subset Z_{H}$-torsor on $X$ we denote by $\Q \cdot \mc{U}$ the $H$-torsor $\Q \times^{Z} \mc{U}$.

 \begin{lem}\label{lem:curlyT}
 Let $T$ be any $F$-rational torus. For any $S = \Spa(R,R^{+}) \in \Perfd_{C}$ and $f \in \Hom_{X_{S}}(t_{S}, T_{S})$, the pullback along \eqref{eq:gerbesm} of the $T_{S}$-torsor on $[\bbD_{S}/\tilde{t}_{S}]$ determined (in the sense of Lemma \ref{lem:equivtors}) by the trivial $T_{S}$-torsor with $\tilde{t}_{S}$-action defined by $f$ is canonically isomorphic to 
 \begin{equation*}
 f_{*}(\mc{T}_{S}^{\tn{loc}}) := \mc{T}_{S}^{\loc} \times^{t_{S},f} T_{S}.
 \end{equation*}
 The above object is viewed as a $t_{S}$-torsor on  $[\mc{T}_{S}^{\loc}/\tilde{t}_{S}]$ via the pullback along  $[\mc{T}_{S}^{\loc}/\tilde{t}_{S}] \to \bbD_{S}$. Moreover, the analogous statement holds with $\bbD_{S}$ replaced with $X_{S}$.
 \end{lem}

 \begin{proof}
     We only prove the $\bbD_{S}$-case (they are identical). By functoriality, it suffices to prove the result for $T = t$ (this is only a pro-torus, but for the purposes of this proof it does not matter) and $f = \tn{id}$. This case follows from the fact that one has a canonical identification $[\mc{T}_{S}^{\loc}/t_{S}] \xrightarrow{\sim} \bbD_{S}$ fitting into the commutative diagram
     \[
     \begin{tikzcd}
         \lb \mc{T}_{S}^{\loc}/\tilde{t}_{S} \rb \arrow{r} \arrow{d} & \lb \bbD_{S}/\tilde{t}_{S} \rb \arrow{d} \arrow{r} & \lb \bbD_{S}/t_{S} \rb \\
         \lb \mc{T}_{S}^{\loc}/t_{S} \rb \arrow{r} &\bbD_{S},\arrow{ru} & 
     \end{tikzcd}
     \]
     where the map $\bbD_{S} \to \lb \bbD_{S}/t_{S} \rb$ is the one associated to $\mc{T}_{S}^{\loc}$, the map $\lb \bbD_{S}/\tilde{t}_{S} \rb \to \lb \bbD_{S}/t_{S} \rb$ is associated to the $t_{S}$-torsor on $\lb \bbD_{S}/\tilde{t}_{S} \rb$ from the statement of the lemma for $T = t$ and $f = \tn{id}$, and the map $[\mc{T}_{S}^{\loc}/\tilde{t}_{S}] \to \bbD_{S}$ induced by the diagram equals the structure map of the gerbe $[\mc{T}_{S}^{\loc}/\tilde{t}_{S}]$.
 \end{proof}

\begin{defi}\label{defi:Qdefs}
For a fixed $S \in \Perfd_{C}$ and $\lambda \in \Hom_{X_{S}}(u_{S}, Z_{G,S})$, we introduce the following objects:
\begin{enumerate}
    \item{Take $\mathfrak{Q}_{\lambda,S}$ to be the $\mathbf{T}_{S}$-torsor on $[\bbD_{S}/\tilde{t}_{S}]$ or $[X_{S}/\tilde{t}_{S}]$ given, via Lemma \ref{lem:equivtors}, by the trivial $\mathbf{T}_{S}$-torsor and the $\tilde{t}_{S}$-action determined by the homomorphism $s(\lambda) \in \Hom_{X_{S}}(\tilde{t}_{S},\mathbf{T}_{S})$. This is where we use the chosen section $s$.}
    \item{Take $\Q_{\lambda,S}'$ to be the pullback of $\mathfrak{Q}_{\lambda,S}$ to $[\mc{T}_{S}^{\loc}/\tilde{t}_{S}]$ or $[\mc{T}_{S}/\tilde{t}_{S}]$ via \eqref{eq:gerbesm} or its global analogue.}
    \item{Take $\widetilde{\Q}_{\lambda,S}$ to be the $G_{S}$-torsor $\mathfrak{Q}_{\lambda,S} \times^{\mathbf{T}_{S}} G_{S}$ corresponding to the embedding $\mathbf{T} \to G$ coming from the choice of pinning.}
\end{enumerate}
\end{defi}

Set $\overline{\mathbf{T}}_{S}:= \mathbf{T}_{S}/Z_{S}$, and, for a $\mathbf{T}_{S}$-torsor $\cP$, denote by $\ov{\cP}$ the $\ov{\mathbf{T}}_{S}$-torsor $\mc{P} \times^{\mathbf{T}_{S}} \ov{\mathbf{T}}_{S}$, similarly with $G_{S}$ and $\ov{G}_{S}$. We alert the reader that we will still frequently use the notation $T/Z$, $G/Z$ despite introducing this notational shortcut---its primary purpose is to lighten the notational load of certain computations.

\begin{lem}\label{lem:Zred}
The trivialization $\psi$ of $\mc{T}_{C}^{\circ}$ from Lemma \ref{lem:Grt} gives a reduction of each $\ov{\mathbf{T}}_{S}$-torsor $\ov{\Q_{\lambda,S}'}|_{[\mc{T}_{S}^{\circ}/\tilde{t}_{S}]}$ and $\ov{\Q_{\lambda,S}'}|_{[\mc{T}_{S}^{\loc,\circ}/\tilde{t}_{S}]}$ to a $Z_{S}$-torsor, which we denote by $\mc{Q}_{\lambda,S}$.
\end{lem}

\begin{proof}
We focus on the $\bbD_{S}$-case. Denoting by $\ov{s(\lambda)} \in \Hom_{X_{S}}(t_{S},\ov{\mathbf{T}}_{S})$ the composition $\tilde{t}_{S} \xrightarrow{s(\lambda)} \mathbf{T}_{S} \to \ov{\mathbf{T}}_{S}$, we observe by Lemma \ref{lem:curlyT} that there is an identification
\begin{equation*}
    \ov{\Q_{\lambda,S}'}  \xrightarrow{\sim} \ov{s(\lambda)}_{*}(\mc{T}_{S}^{\loc}).
\end{equation*}
It follows that, via the pullback of our fixed $X_{C}^{\circ}$-trivialization $\psi$ to $\bbD_{S}^{\circ}$, which we abusively denote simply by $\psi$, we obtain an isomorphism of $\ov{\mathbf{T}}_{S}$-torsors over $\bbD_{S}^{\circ}$
\begin{equation*}
    \ov{\Q_{\lambda,S}'} \xrightarrow{\ov{s(\lambda)}_{*}(\psi)} \underline{\ov{\mathbf{T}}_{S}}.
\end{equation*}
We take our desired reduction $\Q_{\lambda,S}$ to be the fiber over $e \in \underline{\ov{\mathbf{T}}_{S}}(\bbD_{S}^{\circ})$ via the composition
\begin{equation*}
     \Q_{\lambda,S}' \to \ov{\Q_{\lambda,S}'} \xrightarrow{\ov{s(\lambda)}{*}(\psi)} \underline{\ov{\mathbf{T}}_{S}}.
\end{equation*}
The analogue of this construction with $\bbD_{S}$ replaced with $X_{S}$ goes through in an identical manner (since $\psi$ is defined on $X_{C}^{\circ}$).
\end{proof}

\begin{rque}
\begin{enumerate}
    \item{Our notation does not distinguish between the local and global $\Q_{\lambda,S}$-objects because the difference between them causes no risk of confusion---one checks easily that the localization of the global version is canonically isomorphic to the local version.}
    \item{Our above construction gives a map
    \begin{equation*}
        \Hom_{F}(u,Z) \to Z^{1}_{\tn{v}}([\mc{T}_{S}^{\loc,\circ}/\tilde{t}_{S}], Z_{S})
    \end{equation*}
    (one can take \'{e}tale $Z_{S}$-torsors when $F$ has mixed characteristic), which we warn the reader is not multiplicative. This is also true for $\bbD_{S}$ replaced with $X_{S}$.
    }
\end{enumerate}    
\end{rque}

We can finally define the extended Hecke stacks:
\begin{defi}[Extended local Hecke stack with $n$ legs]\label{defi:elochck}
Fix $n \in \mathbb{N}$. Define the functor $\tn{Hck}_{Z \subset G}^{e,\loc,n}$ from $\Perfd_{C}$ to groupoids by sending $S$ to the following data:
\begin{enumerate}
    \item{An $n+1$-tuple $(\cP_{0}, \cP_{1}, \dots, \cP_{n})$ of band-localized (as in Definition \ref{defi:bandloc}) $G_{S}$-torsors on the gerbe $[\mc{T}_{S}^{\loc}/\tilde{t}_{S}]$ such that each $\lambda_{\cP_{i}}$ factors through $Z_{S}$;
    
    \item{A sequence $\{f_{i}\}_{i=0}^{n-1}$, where $f_{i}$ is an isomorphism of $G_{S}$-torsors
    \begin{equation*}
        \Q_{\lambda_{i+1}/\lambda_{i},S} \cdot \cP_{i}|_{[\mc{T}_{S}^{\loc,\circ}/\tilde{t}_{S}]} \xrightarrow{\sim} \cP_{i+1}|_{[\mc{T}_{S}^{\loc,\circ}/\tilde{t}_{S}]},
    \end{equation*}
    where $\lambda_{i}:= \lambda_{\cP_{i}}$.
    }

    \item{A morphism from $((\cP_{i}),(f_{i}))$ to $((\cP_{i}'),(f'_{i}))$ is a sequence $(h_{i})_{i=0}^{n}$ of isomorphisms $\cP_{i} \xrightarrow{h_{i}} \cP_{i}'$ of $G_{S}$-torsors over $[\mc{T}_{S}^{\loc}/\tilde{t}_{S}]$---note that this forces $\lambda_{i} = \lambda_{i}'$---such that $h_{i+1}^{-1} \circ f'_{i} \circ (\Q_{\lambda_{i+1}/\lambda_{i},S} \cdot h_{i}) = f_{i}$ for all $i$.}


 
    }
\end{enumerate}
    
\end{defi}
The global version is completely analogous, but we write it out fully for completeness.

\begin{defi}[Extended global Hecke stack with $n$ legs]\label{defi:eglobhck}
Fix $n \in \mathbb{N}$. Define the functor $\tn{Hck}_{Z \subset G}^{e,n}$ by sending $S$ in $\Perfd_{C}$ to the following data:
\begin{enumerate}
    \item{An $n+1$-tuple $(\cP_{0}, \cP_{1}, \dots, \cP_{n})$ of $G_{S}$-torsors on the gerbe $[\mc{T}_{S}/\tilde{t}_{S}]$ such that each $\lambda_{\cP_{i}}$ factors through $Z_{S}$;}
    
    \item{A sequence $\{f_{i}\}_{i=0}^{n-1}$, where $f_{i}$ is an isomorphism of $G_{S}$-torsors
    \begin{equation*}
        \Q_{\lambda_{i+1}/\lambda_{i},S} \cdot \cP_{i}|_{[\mc{T}_{S}^{\circ}/\tilde{t}_{S}]} \xrightarrow{\sim} \cP_{i+1}|_{[\mc{T}_{S}^{\circ}/\tilde{t}_{S}]} .
    \end{equation*}
    }

 \item{Morphisms in $\tn{Hck}_{Z \subset G}^{e,n}(S)$ are defined in an identical way as for $\tn{Hck}_{Z \subset G}^{e,\loc,n}(S)$.}
\end{enumerate}
\end{defi}

The following is immediate:
\begin{lem}
For $Z \subseteq Z' \subseteq Z_{G}$ there are transition morphisms
\begin{equation}\label{eq:Ztrans}
\tn{Hck}_{Z \subset G}^{e,\loc,n} \to \tn{Hck}_{Z' \subset G}^{e,\loc,n}, \hspace{1mm} \tn{Hck}_{Z \subset G}^{e,n} \to \tn{Hck}_{Z' \subset G}^{e,n}.
\end{equation}
\end{lem}

See Remark \ref{rque:naivereform} for the relationship between the above definitions and Definition \ref{defi:hcknaive} (their ``naive'' analogues).

\begin{rque}\label{rque:non-gluing-to-div-1}
    A notable feature of these extended Hecke stacks is that, for $S = \Spa(R,R^{+}) \in \Perfd_{C}$ every coordinate in $(\Div^{1})^{n}(S)$ is a relative Cartier divisor for $S$ which factors through the map $\Spa(R^{\sharp},R^{\sharp,+}) \xrightarrow{x} \Spa(C^{\sharp},C^{\sharp,+}) \xrightarrow{y} X_{C}$. %In other words, we only allow legs which geometrically factor through the fixed closed point of $X_{C}$ determined by $y$. 
    In rough terms, our constructions of $\tn{Hck}_{Z \subset G}^{e,\loc,n}$ and $\tn{Hck}_{Z \subset G}^{e,n}$ do not ``globalize to all of $X_{C}$''. 
    
    This failure to globalize can be expected, since, as we will see in \S \ref{sec:extended-spectral}, if this construction did globalize (as in the case of the ``naive'' analogues in Definition \ref{defi:hcknaive}), we would obtain a $\tn{Rep}_{\Lambda}(Q^{I})$-linear functor
    \begin{equation*}
        \tn{Rep}_{\Lambda}((\widehat{G/Z} \rtimes Q)^{I}) \to \End_{\Lambda}(\mc{D}_{\tn{lis}}(\Bun_{G}^{e},\Lambda)^{\omega})^{B\Weil_{F}^{I}},
    \end{equation*}
which is unreasonable. In particular, it would imply that every (Fargues--Scholze) $L$-parameter for $\widehat{G}$ can be lifted to an $L$-parameter for $\widehat{G/Z}$.
\end{rque}

\subsection{Convolution on extended Hecke stacks}\label{sec:Hckconv}
We now define a convolution operation on these extended Hecke stacks. The local and global cases are identical, so we construct this in full detail only in the local case. This will be a morphism
\begin{equation*}
    \tn{Hck}_{Z \subset G}^{e,\loc,2} \xrightarrow{\mu} \tn{Hck}_{Z \subset G}^{e,\loc,1}.
\end{equation*}

Let $((\cP_{0},\cP_{1},\cP_{2}), (f_{0},f_{1})) \in \tn{Hck}_{Z \subset G}^{e,\loc,2}(S)$; the goal is to define an isomorphism of $G_{S}$-torsors on $[\mc{T}_{S}^{\loc,\circ}/\tilde{t}_{S}]$
\begin{equation*}
\Q_{\lambda_{2}/\lambda_{0},S} \cdot \cP_{0} \xrightarrow{g} \cP_{2}.
\end{equation*}

We first form the composition
 \begin{equation*}
 \Q_{\lambda_{2}/\lambda_{1},S}\cdot(\Q_{\lambda_{1}/\lambda_{0},S}  \cdot \cP_{0}) \xrightarrow{\Q_{\lambda_{2}/\lambda_{1},S} \cdot f_{0}} \Q_{\lambda_{2}/\lambda_{1},S} \cdot \cP_{1} \xrightarrow{f_{1}} \cP_{2},
 \end{equation*}
 and observe that 
 \begin{equation*}
 \Q_{\lambda_{2}/\lambda_{1},S}\cdot \Q_{\lambda_{1}/\lambda_{0},S} = (\Q_{\lambda_{2}/\lambda_{0},S} \cdot \Q_{\lambda_{2}/\lambda_{0},S}^{-1}) \cdot  \Q_{\lambda_{2}/\lambda_{1},S}\cdot \Q_{\lambda_{1}/\lambda_{0},S} = \Q_{\lambda_{2}/\lambda_{0},S} \cdot [\Q_{\lambda_{2}/\lambda_{0},S}^{-1} \cdot  \Q_{\lambda_{2}/\lambda_{1},S}\cdot \Q_{\lambda_{1}/\lambda_{0},S}]
 \end{equation*}

  To lighten the notational load, temporarily set $a:=z(\lambda_{2}/\lambda_{1},\lambda_{1}/\lambda_{0})$. We note that the right-hand bracketed factor is, by Lemma \ref{lem:curlyT}, canonically identified with the $Z_{S}$-torsor $a_{*}(\mc{T}_{S}^{\loc,\circ})$ on $[\mc{T}_{S}^{\loc,\circ}/\tilde{t}_{S}]$ (viewed as the pullback of the corresponding torsor on $\bbD_{S}^{\circ}$). 
 
    We thus have an isomorphism of $Z_{S}$-torsors
\begin{equation}\label{eq:Arnaudisom}
    \Q_{\lambda_{2}/\lambda_{1},S} \cdot \Q_{\lambda_{1}/\lambda_{0},S} = \Q_{a,S} \cdot \Q_{\lambda_{2}/\lambda_{0},S} = a_{*}(\mc{T}_{S}^{\loc,\circ}) \cdot \Q_{\lambda_{2}/\lambda_{0},S} \xrightarrow{a_{*}(\psi)} \Q_{\lambda_{2}/\lambda_{0},S},
\end{equation}
which we denote by $a_{*}(\psi)$.

We then define the desired isomorphism $g$ to be the composition (over $[\mc{T}_{S}^{\loc,\circ}/\tilde{t}_{S}]$)
 \begin{equation}\label{eq:gconv}
 \Q_{\lambda_{2}/\lambda_{0},S}  \cdot \cP_{0} \xrightarrow{a_{*}(\psi^{-1})}   \Q_{\lambda_{2}/\lambda_{1},S} \cdot (\Q_{\lambda_{1}/\lambda_{0},S} \cdot \cP_{0}) \xrightarrow{\Q_{\lambda_{2}/\lambda_{1},S} \cdot f_{0}} \Q_{\lambda_{2}/\lambda_{1},S} \cdot \cP_{1} \xrightarrow{f_{1}} \cP_{2},
\end{equation}
where we abusively use $a_{*}(\psi^{-1})$ above to denote the isomorphism obtained from the inverse of \eqref{eq:Arnaudisom} by acting on $\cP_{0}$.

 Everything done above works the same with $\bbD_{S}$ and $\bbD_{S}^{\circ}$ replaced with $X_{S}$ and $X_{S}^{\circ}$, defining a convolution morphism
 \begin{equation*}
    \tn{Hck}_{Z \subset G}^{e,2} \xrightarrow{\mu} \tn{Hck}_{Z \subset G}^{e,1}.
\end{equation*}


It follows that for any $1 \leq i \leq n-1$, there are well defined maps
\begin{equation*}
\mu_i : \Hck_{Z \subset G}^{e,\loc, n} \to \Hck_{Z \subset G}^{e,\loc, n-1}, \hspace{1mm} \Hck_{Z \subset G}^{e, n} \to \Hck_{Z \subset G}^{e, n-1}
\end{equation*}
that compose $f_{i-1}$ and $f_{i}$ in the sense defined above. In the case $n=2$, we have $\mu_{1} = \mu$.

\begin{prop}
The above convolution map equips $\Hck_{Z \subset G}^{e,\loc,1}$ with the structure of an algebra in correspondences. More precisely, we have a well-defined functor $\Delta^{\tn{op}} \to \tn{Corr}$ that sends $[1]$ to $\Hck_{Z \subset G}^{e,\loc,1}$ and the map $[1] \to [2]$ to the above multiplication.
\end{prop}

\begin{proof}
This functor is defined by sending $[n]$ to $\Hck^{e,\loc,n}_{Z \subset G}$, and a morphism $[m] \to [n]$ to the corresponding map $\Hck_{Z \subset G}^{e,\loc,n} \to \Hck_{Z \subset G}^{e,\loc,m}$ determined by the convolution defined above. 

The fact that this assignment is unital is clear (by taking the trivial modification), and so it suffices to show that the following associativity diagram commutes
\[
\begin{tikzcd}
    \Hck^{e,\loc,3}_{Z \subset G} \arrow["\mu_{2}"]{d} \arrow["\mu_{1}"]{r} & \Hck^{e,\loc,2}_{Z \subset G} \arrow["\mu"]{d} \\
    \Hck^{e,\loc,2}_{Z \subset G} \arrow["\mu"]{r} & \Hck^{e,\loc,1}_{Z \subset G}.
\end{tikzcd}
\]

For notational brevity, we set $z(\lambda_{i_{1}}/\lambda_{i_{2}},\lambda_{i_{2}}/\lambda_{i_{3}})_{*}(\psi^{-1}) = \psi^{-1}_{i_{1},i_{2},i_{3}}$. First, we observe that the square
\[
\begin{tikzcd}
\Q_{\lambda_{3}/\lambda_{0},S} \cdot \cP_{0} \arrow["\psi^{-1}_{3,2,0}"]{r} \arrow["\psi^{-1}_{3,1,0}"]{d} & \Q_{\lambda_{3}/\lambda_{2},S} \cdot \Q_{\lambda_{2}/\lambda_{0},S} \cdot \cP_{0} \arrow["\Q_{\lambda_{3}/\lambda_{2},S} \cdot \psi^{-1}_{2,1,0}"]{d} \\
\Q_{\lambda_{3}/\lambda_{1},S} \cdot \Q_{\lambda_{1}/\lambda_{0},S} \cdot \cP_{0} \arrow["\psi^{-1}_{3,2,1}"]{r} & \Q_{\lambda_{3}/\lambda_{2},S} \cdot \Q_{\lambda_{2}/\lambda_{1},S} \cdot \Q_{\lambda_{1}/\lambda_{0},S} \cdot \cP_{0}
\end{tikzcd}
\]
commutes, by the two-cocycle condition.

The following square also commutes for elementary reasons:
\[
\begin{tikzcd}[column sep=7em]
    \Q_{\lambda_{3}/\lambda_{1},S} \cdot \Q_{\lambda_{1}/\lambda_{0},S} \cdot \cP_{0} \arrow["\Q_{\lambda_{3}/\lambda_{1},S} \cdot f_{0}"]{r} \arrow["\psi^{-1}_{3,2,1}"]{d} & \Q_{\lambda_{3}/\lambda_{1},S} \cdot \cP_{1} \arrow["\psi^{-1}_{3,2,1}"]{d}\\
    \Q_{\lambda_{3}/\lambda_{2},S} \cdot \Q_{\lambda_{2}/\lambda_{1},S} \cdot \Q_{\lambda_{1}/\lambda_{0},S} \cdot \cP_{0} \arrow["\Q_{\lambda_{3}/\lambda_{2},S} \cdot \Q_{\lambda_{2}/\lambda_{1},S} \cdot f_{0}"]{r} & \Q_{\lambda_{3}/\lambda_{2},S} \cdot \Q_{\lambda_{2}/\lambda_{1},S} \cdot \cP_{1}.
\end{tikzcd}
\]

Starting with $(\cP_{0}, \cP_{1}, \cP_{2}, \cP_{3}, (f_{0},f_{1},f_{2}))$, going right and then going down yields the modification isomorphism
\begin{equation}\label{eq:assoc1}
     f_{2} \circ \Q_{\lambda_{3}/\lambda_{2},S} \cdot [f_{1} \circ (\Q_{\lambda_{2}/\lambda_{1},S} \cdot f_{0}) \circ z(\lambda_{2}/\lambda_{1}, \lambda_{1}/\lambda_{0})_{*}(\psi^{-1})] \circ z(\lambda_{3}/\lambda_{2},\lambda_{2}/\lambda_{0})_{*}(\psi^{-1})
\end{equation}
from $\cP_{0}$ to $\cP_{3}$ (on $[\mathcal{T}_{S}^{\loc,\circ}/\tilde{t}_{S}]$).

Similarly, going the other direction in the diagram yields the modification
\begin{equation}\label{eq:assoc2}
[f_{2} \circ (\Q_{\lambda_{3}/\lambda_{2},S} \cdot f_{1}) \circ z(\lambda_{3}/\lambda_{2},\lambda_{2}/\lambda_{1})_{*}(\psi^{-1})] \circ (\Q_{\lambda_{3}/\lambda_{1},S} \cdot f_{0}) \circ z(\lambda_{3}/\lambda_{1},\lambda_{1}/\lambda_{0})_{*}(\psi^{-1}).
\end{equation}

Noting that the last two maps in both of the above modifications are equal, we are reduced to proving that the corresponding first parts of the compositions are equal (remove the ``$f_{2} \circ \Q_{\lambda_{3}/\lambda_{2},S} \cdot f_{1}$'' terms from the end of both compositions). 

We then deduce the result from the fact that the first part of the modification \eqref{eq:assoc1} is obtained from going right and then down in the first diagram, followed by taking the bottom map in the second diagram, and the first part of the modification \eqref{eq:assoc2} is obtained from going down and then right in the first diagram, followed by taking the long path (up and around) between the bottom two terms in the second diagram. 
\end{proof}

We conclude this subsection by recording some standard compatibilities. 

\begin{prop}\label{prop:hckloc}
    There is a ``global-to-local'' map
    \begin{equation*}
        \tn{Hck}_{Z \subset G}^{e,n} \to \tn{Hck}_{Z \subset G}^{e,\loc,n}
    \end{equation*}
    which for any $i$ makes the following diagram Cartesian
    \[
    \begin{tikzcd}
        \tn{Hck}_{Z \subset G}^{e,n} \arrow["\mu_{i}"]{r} \arrow{d} & \tn{Hck}_{Z \subset G}^{e,n-1} \arrow{d} \\
        \tn{Hck}_{Z \subset G}^{e,\loc,n}  \arrow["\mu_{i}"]{r}  & \tn{Hck}_{Z \subset G}^{
    e,\loc,n-1}.
    \end{tikzcd}
    \]
\end{prop}

\begin{proof}
    This is a straightforward consequence of the fact that our local trivializations $\iota_{\underline{x}}^{*}\psi$ are pulled back from the fixed global trivialization $\psi$.
\end{proof}

\begin{lem}\label{lem:changeinZconv}
   For $Z \subseteq Z' \subseteq Z_{G}$, convolution maps on the local and global extended Hecke stacks are compatible with the transition morphisms \eqref{eq:Ztrans} in the obvious sense.
\end{lem}

\begin{proof}
This is clear.
\end{proof}

\begin{defi}
    We set $\tn{Hck}_{G}^{e,\loc,n} = \varinjlim_{Z} \tn{Hck}_{Z \subset G}^{e,\loc,n}$ and $\tn{Hck}_{G}^{e,n} = \varinjlim_{Z} \tn{Hck}_{Z \subset G}^{e,n}$. 
\end{defi}

By the above two results, we have a localization map $\tn{Hck}_{G}^{e,n} \to \tn{Hck}_{G}^{e,\loc,n} $ and convolution maps
    \begin{equation*}
 \tn{Hck}_{G}^{e,\loc,n} \to  \tn{Hck}_{G}^{e,\loc,n-1},  \hspace{1mm} \tn{Hck}_{G}^{e,n} \to  \tn{Hck}_{G}^{e,n-1}
    \end{equation*}
    which are compatible with localization.



\subsection{Quotient maps}\label{sec:quotmap}
As we will see shortly, a crucial property of the extended Hecke stacks defined above is that they admit maps 
\begin{equation*}
\tn{Hck}^{e,\loc,n}_{Z \subset G} \xrightarrow{\pi} \tn{Hck}^{\loc,n}_{G/Z}, \hspace{1mm} \tn{Hck}^{e,n}_{Z \subset G} \xrightarrow{\pi} \tn{Hck}^{n}_{G/Z}.
\end{equation*}

We explain the construction of the local version of $\pi$; the construction of its global analogue is essentially the same. Recall that for $S \in \Perfd$ we use $\ov{G}_{S}$ to denote $G_{S}/Z_{S}$.

The following general result about descending torsors on gerbes will be used repeatedly:

\begin{lem}\label{lem:descendtors}
    Let $\mathfrak{X} \xrightarrow{f} X$ be a gerbe banded by an abelian group $A$. Let $H$ be a sheaf of groups on $X$, and suppose that $\mc{U}$ is an $f^{*}H$-torsor on $\mathfrak{X}$ such that the inertial action \eqref{eq:bandmap} is trivial. Then $f_{*}\mc{U}$ is an $H$-torsor on $X$ and the adjunction map $f^{*}f_{*}\mc{U} \to \mc{U}$ is an isomorphism of $f^{*}H$-torsors on $\mathfrak{X}$.
\end{lem}

\begin{proof}
    First, observe that we have an identification $f_{*}f^{*}H \xrightarrow{\sim} H$, and so there is an action morphism $f_{*}\mc{U} \times_{X} H \to f_{*}\mc{U}$. It suffices to show that the map $f_{*}\mc{U} \times_{X} H \to f_{*}\mc{U} \times_{X} f_{*}\mc{U}$ is an isomorphism. This may be checked after passing to a cover which trivializes $\mathfrak{X}$, so we can assume that $\mathfrak{X} = [X/A]$.

    For the map $[X/A] \to X$, if we use Lemma \ref{lem:equivtors} to describe $H$-torsors on $[X/A]$, then pushforward sends an $A$-equivariant $H$-torsor $V$ on $[X/A]$ to $V^{A}$. When the $A$-action is trivial, we get the claimed result.
\end{proof}

We will now explain how to map a point $(\cP_{i},(f_{i})) \in \tn{Hck}^{e,\loc,n}_{Z \subset G}(S)$ for $S \in \Perfd_{C}$ to a point of $\Hck_{G/Z}^{\loc,n}(S)$. We define $\pi$ on tuples of torsors by sending $(\cP_{0}, \dots, \cP_{n})$ to $(\ov{\cP_{0}}, \dots, \ov{\cP_{n}})$ (recall also that $\ov{\cP_{i}} := \cP_{i} \times^{G_{S}} \ov{G}_{S}$), which descends canonically (by pushing forward along $[\mc{T}_{S}^{\loc}/\tilde{t}_{S}] \to \bbD_{S}$, using Lemma \ref{lem:descendtors}) to a $\ov{G}_{S}$-torsor on $\bbD_{S}$, also denoted by $\ov{\cP_{i}}$.

Observe that there is a canonical isomorphism 
\begin{equation}\label{eq:varphi}
   \varphi_{\lambda_{i+1}/\lambda_{i},\cP_{i}} \colon \ov{\Q_{\lambda_{i+1}/\lambda_{i},S} \cdot \cP_{i}} \xrightarrow{\sim} \ov{\cP_{i}}
\end{equation}
because, by construction, $\Q_{\lambda_{i+1}/\lambda_{i},S}$ is a $Z_{S}$-torsor.

To define the modifications $(g_{i})$ using the modifications $(f_{i})$ from $\tn{Hck}^{e,\loc,n}_{Z \subset G}(S)$, we take the composition 
\begin{equation*}
g_{i}:= \ov{\cP_{i}} \xrightarrow{\varphi_{\lambda_{i+1}/\lambda_{i},\cP_{i}}^{-1}} \ov{\Q_{\lambda_{i+1}/\lambda_{i},S} \cdot \cP_{i}}  \xrightarrow{\bar{f}_{i}} \ov{\cP_{i+1}},
\end{equation*}
where $\bar{f}_{i} := (-\times^{G_{S}} \ov{G}_{S})(f_{i})$.

Since all of the above constructions are evidently functorial in $S \in \Perfd_{C}$ and globalize, we obtain:
\begin{prop}\label{prop:Hckquot}
    The assignment
    \begin{equation*}
        (\cP_{i},(f_{i})) \mapsto (\ov{\cP_{i}}, (g_{i}))
    \end{equation*}
    defined above gives morphisms of functors over $\Perfd_{C}$
    \begin{equation*}
        \Hck_{Z \subset G}^{e,\loc,n} \xrightarrow{\pi} \Hck_{G/Z}^{\loc,n}, \hspace{1mm} \Hck_{Z \subset G}^{e,n} \xrightarrow{\pi} \Hck_{G/Z}^{n}
    \end{equation*}
    which are compatible with the transition maps \eqref{eq:Ztrans} and the global-to-local maps from Proposition \ref{prop:hckloc}.
\end{prop}

\begin{rque}
    \begin{enumerate}
        \item We will not use the global map $\pi$, but we include it above for completeness.
        \item It is clear that the above map factors canonically as a composition
        \begin{equation*}
             \Hck_{Z \subset G}^{e,\loc,n} \to \Hck_{G/Z,\Spd(C)}^{\loc,n} := \Hck_{G/Z}^{\loc,n} \times_{\Div^{1}} \Spd(C) \to \Hck_{G/Z}^{\loc,n}, 
        \end{equation*}
        and by abuse of notation we will sometimes also use $\pi$ to denote the first map in the composition (this is harmless).
    \end{enumerate}
\end{rque}

The next result is essentially a formal consequence of our construction, but we write out the proof for completeness (the reader is advised to skip it on the first reading).

\begin{prop}\label{prop:Hckconv}
    The following diagram is strictly commutative for any $n$ and $1 \leq i \leq n-1$,
    \[
    \begin{tikzcd}
    \Hck^{e,\loc,n}_{Z \subset G} \arrow["\mu_{i}"]{r} \arrow["\pi"]{d} & \Hck^{e,\loc,n-1}_{Z \subset G}  \arrow["\pi"]{d} \\
    \Hck^{\loc,n}_{G/Z} \arrow["\mu_{i}"]{r} & \Hck^{\loc,n-1}_{G/Z}.
    \end{tikzcd}
    \]
     The same statement holds for the global analogues.
\end{prop}

\begin{proof}
    We prove the local case when $n=2$---the general case follows immediately from this one. 

   For $S \in \Perfd_{C}$, applying the convolution first on a pair  $(\cP_{0}, \cP_{1}, f_{0})$ and $(\cP_{1}, \cP_{2}, f_{1})$ and then applying $\pi$ gives $(\ov{\cP_{0}}, \ov{\cP_{2}}, g)$, where $g$ is the composition 
\begin{equation*}
\ov{\cP_{0}} \xrightarrow{\varphi_{\lambda_{2}/\lambda_{0},\cP_{0}}^{-1}} \ov{\Q_{\lambda_{2}/\lambda_{0},S} \cdot \cP_{0}} \xrightarrow{\ov{z(\lambda_{2}/\lambda_{1},\lambda_{1}/\lambda_{0})_{*}(\psi^{-1})}}  \ov{\Q_{\lambda_{2}/\lambda_{1},S}\cdot(\Q_{\lambda_{1}/\lambda_{0},S}  \cdot \cP_{1})} 
\end{equation*}
followed by 
\begin{equation*}
\ov{\Q_{\lambda_{2}/\lambda_{1},S}\cdot(\Q_{\lambda_{1}/\lambda_{0},S}  \cdot \cP_{1})}  \xrightarrow{\ov{\Q_{\lambda_{2}/\lambda_{1},S} \cdot f_{0}}} \ov{\Q_{\lambda_{2}/\lambda_{1},S} \cdot \cP_{1}} \xrightarrow{\bar{f_{1}}} \ov{\cP_{2}}.
\end{equation*}

On the other hand, applying $\pi$ and then convolving gives the composition
\begin{equation*}
\ov{\cP_{0}} \xrightarrow{\varphi_{\lambda_{1}/\lambda_{0},\cP_{0}}^{-1}} \ov{\Q_{\lambda_{1}/\lambda_{0},S}\cdot \cP_{0}}  \xrightarrow{\bar{f}_{0}} \ov{\cP_{1}} \xrightarrow{\varphi_{\lambda_{2}/\lambda_{1},\cP_{1}}^{-1}} 
\ov{\Q_{\lambda_{2}/\lambda_{1},S}\cdot \cP_{1}}  \xrightarrow{\bar{f}_{1}} \ov{\cP_{2}}.
\end{equation*}
Equality of the two compositions then reduces to proving the equality
\begin{equation*}
\ov{\Q_{\lambda_{2}/\lambda_{1},S} \cdot f_{0}} \circ  \ov{z(\lambda_{2}/\lambda_{1},\lambda_{1}/\lambda_{0})_{*}(\psi^{-1})} \circ \varphi_{\lambda_{2}/\lambda_{0},\cP_{0}}^{-1} = \varphi_{\lambda_{2}/\lambda_{1},\cP_{1}}^{-1} \circ \bar{f}_{0} \circ \varphi_{\lambda_{1}/\lambda_{0},\cP_{0}}^{-1},
\end{equation*}
and therefore further, using $\varphi_{\lambda_{2}/\lambda_{1},\cP_{1}}^{-1} \circ \bar{f}_{0} = \ov{\Q_{\lambda_{2}/\lambda_{1},S} \cdot f_{0}} \circ \varphi_{\lambda_{2}/\lambda_{1},\cP_{1}}^{-1}$, to the equality 
\begin{equation}\label{eq:varphieq1}
    \ov{z(\lambda_{2}/\lambda_{1},\lambda_{1}/\lambda_{0})_{*}(\psi^{-1})} \circ \varphi_{\lambda_{2}/\lambda_{0},\cP_{0}}^{-1} =  \varphi_{\lambda_{2}/\lambda_{1},\cP_{1}}^{-1} \circ \varphi_{\lambda_{1}/\lambda_{0},\cP_{0}}^{-1}.
\end{equation}
We note that, by definition,
\begin{equation*}
\ov{z(\lambda_{2}/\lambda_{1},\lambda_{1}/\lambda_{0})_{*}(\psi^{-1})} = \varphi^{-1}_{z(\lambda_{2}/\lambda_{1},\lambda_{1}/\lambda_{0}),\cP_{0}}
\end{equation*}
and that, as $Z_{S}$-torsors,
\begin{equation*}
    \Q_{z(\lambda_{2}/\lambda_{1},\lambda_{1}/\lambda_{0}),S} \cdot \Q_{\lambda_{2}/\lambda_{0},S}= \Q_{\lambda_{2}/\lambda_{1},S} \cdot \Q_{\lambda_{1}/\lambda_{0},S} ,
\end{equation*}
and so we deduce the equality \eqref{eq:varphieq1}.
\end{proof}

Recall that we have a locally constant morphism
\begin{equation*}
    \Hck^{e,\loc,1}_{Z \subset G} \xrightarrow{\nu^{e}} \underline{\Hom}_{X}(u,Z)^{2}
\end{equation*}
(where $\underline{\Hom}_{X}(u,Z)$ is from Definition \ref{defi:hom}) given by sending $(\cP_{0}, \cP_{1}, f)$ to $(\lambda_{0},\lambda_{1})$, which induces a decomposition (as v-stacks)
\begin{equation}\label{eq:HckeDecomp}
\Hck^{e,\loc,1}_{Z \subset G} = \bigsqcup_{(\lambda_{0},\lambda_{1}) \in \Hom_{F}(u,Z)^{2}} \Hck^{e,\loc,1,(\lambda_{0},\lambda_{1})}_{Z \subset G}.
\end{equation}

\begin{rque}\label{rque:naivereform}
It is easy to see that 
\begin{equation*}
\Hck^{e,\loc,\{0,1\},\tn{naive}}_{Z \subset G} \times_{\Div^{1}} \Spd(C) = \bigsqcup_{\lambda \in \Hom_{F}(u,Z)} \Hck_{Z \subset G}^{e,\loc,1,(\lambda,\lambda)},
\end{equation*}
similarly with the global versions.
\end{rque}

There is an analogous decomposition for $\Hck_{G/Z}^{\loc,1}$. Recall that, by \cite[Proposition 21.1.4]{SW20}, $\pi_{0}(\Hck^{\loc,1}_{G/Z})$ is canonically identified with $\pi_{1}(G/Z)$, which fits into an exact sequence
\begin{equation*}
    0 \to \pi_{1}(G) \to \pi_{1}(G/Z) \to \Hom_{\ov{F}}(\mu,Z) \to 0.
\end{equation*}
We thus obtain a locally constant map
\begin{equation*}
\Hck^{\loc,1}_{G/Z} \xrightarrow{\nu} \Hom_{\ov{F}}(\mu, Z)
\end{equation*}
by taking the composition
\begin{equation*}
\Hck^{\loc,1}_{G/Z} \to \pi_{0}(\Hck^{\loc,1}_{G/Z}) \to \Hom_{\ov{F}}(\mu,Z),
\end{equation*}
as well as a decomposition (as v-stacks)
\begin{equation*}
\Hck^{\loc,1}_{G/Z} = \bigsqcup_{\lambda \in \Hom_{\ov{F}}(\mu, Z)} \Hck^{\loc,1,\lambda}_{G/Z}.
\end{equation*}
One defines $\Hck^{\loc,2,(\nu_{0},\nu_{1})}_{G/Z}$ for $\nu_{i} \in \Hom_{\ov{F}}(\mu,Z)$ by mapping $\Hck^{\loc,2}_{G/Z}$ to $\Hck^{\loc,1}_{G/Z} \times \Hck^{\loc,1}_{G/Z}$ via $(\mc{P}_{0},\mc{P}_{1},\mc{P}_{2},(f_{i})) \mapsto (\mc{P}_{0},\mc{P}_{1},f_{0}) \times (\mc{P}_{1},\mc{P}_{2},f_{1})$ and then applying $\nu \times \nu$. 

We make the elementary but important observation that we have an isomorphism
\begin{equation}\label{eq:utomu}
    \Hom_{F}(u,Z) \xrightarrow{\sim} \Hom_{\ov{F}}(\mu,Z)
\end{equation}
given by $\lambda \mapsto \lambda \circ \widetilde{\tn{Sh}}$, where $\widetilde{\tn{Sh}}$ is the map $\mu_{\ov{F}} \to u_{\ov{F}}$ given by the morphism on character modules 
\begin{equation*}
    \mathcal{C}(\Gal_{F}, \mathbb{Q}/\Z) \to \mathbb{Q}/\Z
\end{equation*}
sending $f$ to $f(e)$.

\begin{prop}\label{prop:Quotslope}
  For $\lambda_{0}, \lambda_{1} \in \Hom_{F}(u,Z)$, the substack $\Hck_{Z \subset G}^{e,\loc,1,(\lambda_{0},\lambda_{1})}$ maps via $\pi$ into $\Hck_{G/Z}^{\loc,1,(\lambda_{0}/\lambda_{1}) \circ \widetilde{\tn{Sh}}}$.  
\end{prop}

\begin{proof}
The fact that $\nu^{e}$ is locally constant implies that it suffices to show that $\Hck_{Z \subset G}^{e,\loc,1,(\lambda_{0},\lambda_{1})}(S)$ maps into $\Hck_{G/Z}^{\loc,1,(\lambda_{0}/\lambda_{1}) \circ \widetilde{\tn{Sh}}}(S)$ for $S = \Spa(R,R^{+}) \in \Perfd_{C}$ with $R$ an algebraically closed field. By the proof of Lemma \ref{lem:cohomvanish} (with $t$ replaced by $G$), we may assume that the fixed $X \in \Hck_{Z \subset G}^{e,\loc,1,(\lambda_{0},\lambda_{1})}(S)$ is of the form $(\widetilde{\Q}_{\lambda_{0},S}, \widetilde{\Q}_{\lambda_{1},S},f)$ (see Definition \ref{defi:Qdefs} to recall the definition of $\widetilde{\Q}_{\lambda_{i},S}$). 


Moreover, for a $Z_{S}$-torsor $\mc{U}$ on  $[\mc{T}_{S}^{\loc}/\tilde{t}_{S}]$, the tuple  $(\mc{U} \cdot \widetilde{\Q}_{\lambda_{0},S}, \mc{U} \cdot \widetilde{\Q}_{\lambda_{1},S},\mc{U} \cdot f)$ also lies in $\Hck_{Z \subset G}^{e,\loc,1}(S)$ and has the same image in $\Hck_{G/Z}^{\loc,1}(S)$ as $(\widetilde{\Q}_{\lambda_{0},S}, \widetilde{\Q}_{\lambda_{1},S},f)$. This lets us replace $X$ by a $\mc{U}$-translate to further assume that $X = (\underline{G_{S}},\widetilde{\Q}_{\lambda,S},f)$ for some $\lambda \in \Hom_{F}(u,Z)$. Finally, it is straightforward to check that replacing $f$ in $X$ with $f \circ (\Q_{\lambda,S} \cdot a)$ for any automorphism of torsors $\underline{G_{S}}|_{\bbD_{S}} \xrightarrow{a} \underline{G_{S}}|_{\bbD_{S}}$ does not affect the image of $\pi(X)$ in
$\Hom_{F}(u,Z) = \pi_{1}(G/Z)/\pi_{1}(G)$. 

Recall the isomorphism 
\begin{equation*}
    \Q_{\lambda,S} \cdot \underline{G_{S}} \xrightarrow{h} \widetilde{\Q}_{\lambda,S},
\end{equation*}
where $h$ is the canonical map 
\begin{equation*}
     \Q_{\lambda,S} \times^{Z_{S}} \underline{G_{S}} \xrightarrow{h'} \mathfrak{Q}_{\lambda,S} \times^{\mathbf{T}_{S}} \underline{G_{S}} = \widetilde{\Q}_{\lambda,S}
\end{equation*}
(again, see Definition \ref{defi:Qdefs} to recall $\mathfrak{Q}_{\lambda,S}$---here $h'$ is obtained via the isomorphism from $\Q_{\lambda,S} \times^{Z_{S}} \mathbf{T}_{S}$ to $\mathfrak{Q}_{\lambda,S}$ implicit in the definition of $\Q_{\lambda,S}$). Since any given isomorphism
\begin{equation*}
    \Q_{\lambda,S} \cdot \underline{G_{S}} \xrightarrow{f} \widetilde{\Q}_{\lambda,S}
\end{equation*}
differs from the map $h$ by such an automorphism $a$, it is enough to prove the result for the object $X =(\underline{G_{S}},\widetilde{\Q}_{\lambda,S},h) \in \Hck_{Z \subset G}^{e,\loc,1,(1,\lambda)}(S)$. By construction, $X$ is the image of the element $(\underline{\mathbf{T}_{S}}, \mathfrak{Q}_{\lambda,S},h') \in \Hck_{Z \subset \mathbf{T}}^{e,\loc,1,(1,\lambda)}(S)$. Since we have a commutative diagram
\[
\begin{tikzcd}
\Hck^{e,\loc,1}_{Z \subset \mathbf{T}}(S) \arrow{r} \arrow{d} & \Hck_{\mathbf{T}/Z}^{\loc,1}(S) \arrow{r} \arrow{d} & X_{*}(\mathbf{T}/Z) \arrow{r} \arrow{d} & \Hom_{F}(u,Z) \arrow["\tn{id}"]{d} \\
\Hck_{Z \subset G}^{e,\loc,1}(S) \arrow{r} & \Hck_{G/Z}^{\loc,1}(S) \arrow{r}  &  \pi_{1}(G/Z)\arrow{r} & \Hom_{F}(u,Z),
\end{tikzcd}
\]
we can reduce further to the case where $G=\mathbf{T}$ is a torus.

By definition, the image $X'$ of $X$ in $\Hck_{\mathbf{T}/Z}^{\loc,1}(S)$ is $(\underline{\ov{\mathbf{T}}_{S}},\ov{\mathfrak{Q}_{\lambda,S}},\bar{h} \circ \varphi_{\lambda,\underline{\ov{\mathbf{T}}_{S}}}^{-1})$---after switching the two torsors, which negates the corresponding point of $\pi_{0}(\Hck_{\mathbf{T}/Z}^{\loc,1})$ by \cite[\S I.6]{Geometrization}, we can and do view this as a point in the affine Grassmannian $\tn{Gr}_{\mathbf{T}/Z}(S)$. By the way we have set things up, this resulting point in $\tn{Gr}_{\mathbf{T}/Z}(S)$ is the image of the point $(\underline{\ov{\mathbf{T}}_{S}},\mc{T}_{S}^{\loc},\psi^{-1}) \in \tn{Gr}_{t}(S)$ via the map $t \xrightarrow{\ov{s(\lambda)}} \mathbf{T}/Z$. We deduce from the commutativity of  
\[
\begin{tikzcd}
\tn{Gr}_{t}(S) \arrow{d} \arrow["\ov{s(\lambda)}"]{r} & \tn{Gr}_{\mathbf{T}/Z}(S) \arrow{d} \\
X_{*}(t) \arrow["\ov{s(\lambda)} \circ -"]{r} & X_{*}(\mathbf{T}/Z)
\end{tikzcd}
\]
(and the definition of $\psi$ from Lemma \ref{lem:Grt}) that the image of $X'$ in $X_{*}(\mathbf{T}/Z)$ is $\ov{s(\lambda)} \circ \tn{Sh}$, whose image in $\Hom_{\ov{F}}(\mu,Z)$ is exactly $\lambda \circ \widetilde{\tn{Sh}}$. Applying the inversion from switching the torsors yields $(1/\lambda) \circ \widetilde{\tn{Sh}}$, as claimed.
\end{proof}

Henceforth, we will implicitly use the isomorphism \eqref{eq:utomu} to identify $\Hom_{F}(u,Z)$ with $\Hom_{\ov{F}}(\mu,Z)$ rather than writing the map explicitly. For example, we write $\Hck_{G/Z}^{\loc,1,\lambda_{0}/\lambda_{1}}$ to denote the stack $\Hck_{G/Z}^{\loc,1,(\lambda_{0}/\lambda_{1}) \circ \widetilde{\tn{Sh}}}$.

\begin{lem}\label{lem:outer-faces}
Fix $\lambda_0,\lambda_1,\lambda_2\in\Hom_F(u,Z)$. The diagram
\[
\begin{tikzcd}[column sep=huge, row sep=large]
\Hck^{e,\mathrm{loc},1,(\lambda_0,\lambda_1)}_{Z\subset G}
  \arrow[d,"\pi"']
& \Hck^{e,\mathrm{loc},2,(\lambda_0,\lambda_1,\lambda_2)}_{Z\subset G}
  \arrow[l]
  \arrow[r]
  \arrow[d,"\pi"]
& \Hck^{e,\mathrm{loc},1,(\lambda_1,\lambda_2)}_{Z\subset G}
  \arrow[d,"\pi"]
\\
\Hck^{\mathrm{loc},1,\lambda_0/\lambda_1}_{G/Z}
& \Hck^{\mathrm{loc},2,(\lambda_0/\lambda_1,\,\lambda_1/\lambda_2)}_{G/Z}
  \arrow[l]
  \arrow[r]
& \Hck^{\mathrm{loc},1,\lambda_1/\lambda_2}_{G/Z}
\end{tikzcd}
\]
is strictly commutative, where the vertical maps are those of Proposition \ref{prop:Hckquot}. In particular, $\pi$ carries $\Hck^{e,\mathrm{loc},2,(\lambda_0,\lambda_1,\lambda_2)}_{Z\subset G}$ into $\Hck^{\mathrm{loc},2,(\lambda_0/\lambda_1,\,\lambda_1/\lambda_2)}_{G/Z}$.
\end{lem}

\begin{proof}
    This is immediate.
\end{proof}

\subsection{Extended affine Grassmannians}
To understand perverse sheaves on $\Hck^{e,\loc,1}_{Z \subset G}$ we need to introduce and study a version of the affine Grassmannian which is adapted to the extended Hecke stack $\Hck^{e,\loc,1}_{Z \subset G}$.

First, denote the local Hecke correspondence by
\[
\begin{tikzcd}
    & \Hck_{Z \subset G}^{e,\loc,1} \arrow["\overleftarrow{h}"]{ld} \arrow["\overrightarrow{h}"]{rd}& \\
    B^{e}L^{+}G & & B^{e}L^{+}G,
\end{tikzcd}
\]
where $B^{e}L^{+}G$ denotes the functor on $\Perfd_{C}$ sending $S$ to the groupoid of $G_{S}$-torsors on $[\mc{T}_{S}^{\loc}/\tilde{t}_{S}]$.


\begin{defi}\label{defi:Gre} \begin{enumerate}
   \item{ Define the functor $\tn{Gr}_{Z \subset G}^{e,(\lambda_{0},\lambda_{1})}$ from $\Perfd_{C}$ to groupoids as the fiber of $\widetilde{\Q}_{\lambda_{1},-}$ under $\overrightarrow{h}|_{\Hck_{Z \subset G}^{e,\loc,1,(\lambda_{0},\lambda_{1})}}$. 
   
   Explicitly, this is the functor sending $S$ to all pairs $(\cP, f)$, where $\cP$ is a band-localized $G_{S}$-torsor on $[\mc{T}_{S}^{\loc}/\tilde{t}_{S}]$ whose inertial morphism is $\lambda_{0}$ and $f$ is an isomorphism of  $G_{S}$-torsors over $[\mc{T}_{S}^{\loc,\circ}/\tilde{t}_{S}]$
     \begin{equation*}
        \Q_{\lambda_{1}/\lambda_{0},S} \cdot \cP \xrightarrow{f} \widetilde{\Q}_{\lambda_{1},S}.
     \end{equation*}
     Morphisms between $(\cP,f)$ and $(\cP',f')$ are isomorphisms of $G_{S}$-torsors over $[\mc{T}_{S}^{\loc}/\tilde{t}_{S}]$ which are compatible with $f$ and $f'$. }

    \item{More generally, we can define $\tn{Gr}_{Z \subset G}^{e,\lambda}$ as the fiber of $\widetilde{\Q}_{\lambda,-}$ under $\overrightarrow{h}$.
         It is easy to see that we have a decomposition
     \begin{equation*}
         \tn{Gr}_{Z \subset G}^{e,\lambda} = \bigsqcup_{\lambda' \in \Hom_{F}(u,Z)} \tn{Gr}_{Z \subset G}^{e,(\lambda',\lambda)}.
     \end{equation*}}

\item{We set $\tn{Gr}_{Z \subset G}^{e} = \bigsqcup_{\lambda \in \Hom_{F}(u,Z)} \tn{Gr}_{Z \subset G}^{e,\lambda}$.}
     \end{enumerate}
     
\end{defi}

By analogy to the usual affine Grassmannian, we have:
\begin{lem}\label{lem:Grrig}
    The automorphism group of a fixed point $X \in \tn{Gr}_{Z \subset G}^{e,\lambda}(S)$ for $S \in \Perfd_{C}$ is trivial.
\end{lem}

\begin{proof}
    Let $a$ be an automorphism of the $G_{S}$-torsor $\cP$ over $[\mc{T}_{S}^{\loc}/\tilde{t}_{S}]$ such that $f \circ (\Q_{\lambda,S} \cdot a)  = f$. It follows that  $\Q_{\lambda,S} \cdot a$ is trivial as an automorphism of $\Q_{\lambda,S} \cdot \cP$ over $[\mc{T}_{S}^{\loc,\circ}/\tilde{t}_{S}]$, and is therefore trivial on $[\mc{T}_{S}^{\loc}/\tilde{t}_{S}]$.
\end{proof}

Since $\pi_{0}(\tn{Gr}_{G}) = \pi_{0}(\Hck_{G}^{\loc,1}) = \pi_{1}(G)$ for any connected reductive $G$, we obtain, as in the previous subsection, a locally constant map
\begin{equation*}
\nu \colon \tn{Gr}_{G/Z,\Spd(C)} \to \pi_{1}(G/Z) \to \Hom_{\ov{F}}(\mu,Z) = \Hom_{F}(u,Z),
\end{equation*}
and a decomposition 
\begin{equation*}
  \tn{Gr}_{G/Z,\Spd(C)}   = \bigsqcup_{\lambda \in \Hom_{F}(u,Z)} \tn{Gr}_{G/Z,\Spd(C)}^{\lambda}.
\end{equation*}

\begin{prop}\label{prop:Grquot}
    There is a functor between groupoids over $\Perfd_{C}$
    \begin{equation*}
\tn{Gr}_{Z \subset G}^{e,(\lambda_{0},\lambda_{1})} \xrightarrow{\pi^{(\lambda_{0},\lambda_{1})}} \tn{Gr}_{G/Z,\Spd(C)}^{\lambda_{0}/\lambda_{1}}
    \end{equation*}
    making the following diagram commute up to $2$-isomorphism:
\[
\begin{tikzcd}
    \Hck_{Z \subset G}^{e,\loc,1,(\lambda_{0},\lambda_{1})} \arrow["\pi"]{r} & \Hck_{G/Z}^{\loc,1,\lambda_{0}/\lambda_{1}} \\
    \tn{Gr}_{Z \subset G}^{e,(\lambda_{0},\lambda_{1})} \arrow{u} \arrow["\pi^{(\lambda_{0},\lambda_{1})}"]{r} & \tn{Gr}_{G/Z,\Spd(C)}^{\lambda_{0}/\lambda_{1}} \arrow{u},
\end{tikzcd}
    \]
    where the top map is from Proposition \ref{prop:Hckquot} (and we are also implicitly using Proposition \ref{prop:Quotslope}).
\end{prop}

\begin{proof}
Recall the fixed auxiliary $\bbD_{C}$-trivialization $\tilde{\psi}$ of $\mc{T}_{C}^{\loc}$ from \S \ref{sec:affGrasscalc}. Given an $S$-point $(\cP, f)$ on the left-hand side, we post-compose the map 
   \begin{equation*}
   \ov{\cP} \xrightarrow{\varphi_{\lambda_{1}/\lambda_{0},\cP}} \ov{\Q_{\lambda_{1}/\lambda_{0},S} \cdot \cP} \xrightarrow{\bar{f}} \ov{\widetilde{\Q}_{\lambda_{1},S}}
   \end{equation*}
   with the trivialization
   \begin{equation*}
  \ov{\widetilde{\Q}_{\lambda_{1},S}} = \ov{s(\lambda_{1})}_{*}(\mc{T}^{\loc}_{S}) \xrightarrow{\ov{s(\lambda_{1})}_{*}(\tilde{\psi})} \underline{\ov{G}_{S}}
   \end{equation*}
   to obtain the point $(\ov{\cP}, \ov{s(\lambda_{1})}_{*}(\tilde{\psi}) \circ \bar{f} \circ \varphi_{\lambda_{1}/\lambda_{0},\cP}) \in \tn{Gr}_{G/Z,\Spd(C)}(S)$, which we take to be the image of our object under $\pi^{(\lambda_{0},\lambda_{1})}$. It is straightforward to check that a different choice of $\tilde{\psi}$ as above gives an isomorphic object in $\tn{Gr}_{G/Z,\Spd(C)}(S)$ (in a way which is functorial in $S$). An isomorphism $h$ simply goes to $\bar{h}$, finishing the definition of $\pi^{(\lambda_{0},\lambda_{1})}$.
   
     The fact that the above map lands in $\tn{Gr}_{G/Z,\Spd(C)}^{\lambda_{0}/\lambda_{1}}$ and commutes up to $2$-isomorphism is immediate from the definition of the (extended and classical) affine Grassmannians and the compatibility of the two local Hecke correspondences.
\end{proof}

\begin{lem}\label{lem:Grtwist} There is an isomorphism 
\begin{equation*}
    \tn{Gr}^{e,(\lambda_{0},\lambda_{1})}_{Z \subset G} \xrightarrow{\sim} \tn{Gr}^{e,(1,\lambda_{1}/\lambda_{0})}_{Z \subset G}
\end{equation*}
over $\tn{Gr}_{G/Z,\Spd(C)}$, defined by sending $(\cP, f)$ to $(\Q_{\lambda_{0}^{-1},S} \cdot \cP, z(\lambda_{1},\lambda_{0}^{-1})_{*}(\psi^{-1}) \circ (\Q_{\lambda_{0}^{-1},S} \cdot f))$. 
\end{lem}
\begin{proof}
    This is a straightforward computation.
\end{proof}

\begin{prop}\label{prop:mainGrisom}
    The map $\pi^{(\lambda_{0},\lambda_{1})}$ is an equivalence of $\tn{v}$-stacks. 
\end{prop}

\begin{proof}
Lemma \ref{lem:Grtwist} reduces us to the case where $\lambda_{0}=1$ and we set $\lambda_{1}/\lambda_{0} = \lambda$, which simplifies computations. Next, observe that there is an isomorphism of functors (which is not $L^{+}G$-equivariant) over $\Perfd_{C}$
    \begin{equation}\label{eq:Grpfeq1}
        \tn{Gr}_{G,\Spd(C)} \xrightarrow{\sim} \tn{Gr}_{Z \subset G}^{e,(1,\lambda)}
    \end{equation}
    defined by sending $(\cP,f)$ to the torsor $\cP$ along with the isomorphism
\begin{equation*}
    \Q_{\lambda,S} \cdot \cP \xrightarrow{\Q_{\lambda,S}\cdot f} \Q_{\lambda,S} \cdot \underline{G_{S}} = \widetilde{\Q}_{\lambda,S}.
\end{equation*}

On the other hand, one has a (non-$L^{+}G$-equivariant) ``translation'' isomorphism 
 \begin{equation}\label{eq:Grpfeq2}
     \tn{Gr}_{G/Z,\Spd(C)}^{1} \xrightarrow{\sim} \tn{Gr}_{G/Z,\Spd(C)}^{1/\lambda}
 \end{equation}
 given by sending $(\mc{M},f)$ to $(\mc{M}, \ov{s(\lambda)}_{*}(\tilde{\psi}) \circ \ov{s(\lambda)}_{*}(\psi^{-1}) \circ f)$, cf. the proof of Proposition \ref{prop:Quotslope} (and $\tilde{\psi}$ is the auxiliary $\bbD_{C}$-trivialization from the end of \S \ref{sec:affGrasscalc}). The map $\tn{Gr}_{G,\Spd(C)} \to \tn{Gr}_{G/Z,\Spd(C)}$ factors as the composition
 \begin{equation*}
     \tn{Gr}_{G,\Spd(C)} \to \tn{Gr}_{G/Z,\Spd(C)}^{1} \hookrightarrow \tn{Gr}_{G/Z,\Spd(C)}.
 \end{equation*}

 Now, observing that the following diagram commutes
 \[
\begin{tikzcd}
\tn{Gr}_{G,\Spd(C)} \arrow["\eqref{eq:Grpfeq1}"]{r} \arrow["\tn{id}"]{d} & \tn{Gr}_{Z \subset G}^{e,(1,\lambda)} \arrow["\pi^{(1,\lambda)}"]{r} & \tn{Gr}_{G/Z,\Spd(C)}^{1/\lambda} \arrow["\tn{id}"]{d} \\
\tn{Gr}_{G,\Spd(C)} \arrow{r} & \tn{Gr}_{G/Z,\Spd(C)}^{1} \arrow["\eqref{eq:Grpfeq2}"]{r} & \tn{Gr}_{G/Z,\Spd(C)}^{1/\lambda},
\end{tikzcd}
 \]
 we deduce the result from the fact that the bottom line is an equivalence.
\end{proof}

\begin{corol}\label{rque:sheafequiv}
   Pullback along $\pi^{(\lambda_{0},\lambda_{1})}$ induces an isomorphism between the categories of sheaves of abelian groups on $\tn{Gr}_{Z \subset G}^{e,(\lambda_{0},\lambda_{1})}$ and $\tn{Gr}_{G/Z,\Spd(C)}^{\lambda_{0}/\lambda_{1}}$.  
\end{corol}

\begin{proof}
     This comes from combining Proposition \ref{prop:mainGrisom} with Lemma \ref{lem:Grrig}.
\end{proof}

\begin{defi} We denote by $\mc{A}_{\lambda_{1}}$ the functor on $\Perfd_{C}$ sending $S$ to $\underline{\tn{Aut}}_{[\mc{T}_{S}^{\loc}/\tilde{t}_{S}]}(\widetilde{\Q}_{\lambda_{1},S})$. 
\end{defi}


\begin{rque}\label{rque:Alambda} \begin{enumerate}
    \item Using the homomorphism $\mc{A}_{\lambda_{1}} \to L^{+}(G/Z)$ induced by $\tilde{\psi}$, the map $\pi^{(\lambda_{0},\lambda_{1})}$ is $\mc{A}_{\lambda_{1}}$-equivariant. 

    \item Note that $\tilde{\psi}$ induces an isomorphism of gerbes $[\mc{T}_{C}^{\loc}/\tilde{t}_{C}] \xrightarrow{\sim} [\bbD_{C}/u_{C}]$ over $\bbD_{C}$ under which $\widetilde{\Q}_{\lambda_{1},C}$ corresponds to the trivial $G_{C}$-torsor with $u_{C}$-action given by $\lambda_{1}$ (using Lemma \ref{lem:equivtors}). Since $\lambda_{1}$ is valued in $Z_{G}$, every $\bbD_{C}$-automorphism of the trivial $G_{C}$-torsor is $u_{C}$-equivariant, giving an isomorphism $\mc{A}_{\lambda_{1}} \xrightarrow{\sim} L^{+}G$ which is compatible with the maps from both groups to $L^{+}(G/Z)$.

\end{enumerate}
\end{rque}


\begin{rque}\label{rque:strat}
    Proposition \ref{prop:mainGrisom} shows that, via the isomorphism $\pi^{(\lambda_{0},\lambda_{1})}$, the locally spatial diamond $\tn{Gr}_{Z \subset G}^{e,(\lambda_{0},\lambda_{1})}$ inherits a natural stratification from the Schubert stratification of $\tn{Gr}_{G/Z,\Spd(C)}^{\lambda_{0}/\lambda_{1}}$ which is indexed by all $\nu \in X_{*}(T/Z)^{+,\lambda_{0}/\lambda_{1}}$, where the superscript ``$\lambda_{0}/\lambda_{1}$'' means that the image of $\nu$ in $\Hom_{\ov{F}}(\mu, Z)$ is $\lambda_{0}/\lambda_{1}$. The same then holds for $\Hck_{Z \subset G}^{e,\loc,1}$. In particular, we can talk about the Schubert cells $\tn{Gr}_{Z \subset G,\nu}^{e,(\lambda_{0},\lambda_{1})}$ and $\Hck_{Z \subset G,\nu}^{e,\loc,1}$ for $\nu \in X_{*}(T/Z)^{+,\lambda_{0}/\lambda_{1}}$.
\end{rque}

By construction, there is an identification
\begin{equation*}
 \Hck_{Z \subset G}^{e,\loc,1,(\lambda_{0},\lambda_{1})} =   \mc{A}_{\lambda_{1}} \backslash \tn{Gr}_{Z \subset G}^{e,(\lambda_{0},\lambda_{1})}.
\end{equation*}

\begin{proof}
This is a consequence of the equalities $\Hck_{Z \subset G}^{e,\loc,1,(\lambda_{0},\lambda_{1})} =  \mc{A}_{\lambda_{1}} \backslash \tn{Gr}_{Z \subset G}^{e,(\lambda_{0},\lambda_{1})}$, $\Hck_{G/Z}^{\loc,1,\lambda_{0}/\lambda_{1}} = L^{+}(G/Z) \backslash \tn{Gr}_{G/Z,\Spd(C)}^{\lambda_{0}/\lambda_{1}}$, Proposition \ref{prop:mainGrisom}, and the fact that the kernel of $\mc{A}_{\lambda_{1}} \to L^{+}(G/Z)$ is $Z$.
\end{proof}

Finally, all of the constructions in this section are functorial in $Z \subseteq Z' \subseteq Z_{G}$, in the sense that there are ``transition morphisms''
\begin{equation*}
\tn{Gr}_{Z \subset G}^{e} \to \tn{Gr}_{Z' \subset G}^{e}
\end{equation*}
which are compatible with the map from Proposition \ref{prop:Grquot} and the above identifications with $\Hck_{Z \subset G}^{e,\loc,1}$ and  $\Hck_{Z' \subset G}^{e,\loc,1}$.

\subsection{Geometric Satake}
 We assume that our coefficient ring $\Lambda$  is killed by $n$ coprime to $p$.
 
 \begin{defi}\label{defi:convGre} Given a triple $\lambda_{0},\lambda_{1}, \lambda_{2} \in \Hom_{F}(u,Z)$, define the \textbf{extended convolution affine Grassmannian} $\tn{Gr}_{Z \subset G}^{e,(\lambda_{0},\lambda_{1},\lambda_{2})}$ to be the fiber over $\widetilde{\Q}_{\lambda_{2},-}$ under $\overrightarrow{h}|_{\Hck_{Z \subset G}^{e,\loc,2,(\lambda_{0},\lambda_{1},\lambda_{2})}}$.
 \end{defi}

The convolution morphism from \S \ref{sec:Hckconv} induces a convolution morphism
\begin{equation*}
    \tn{Gr}_{Z \subset G}^{e,(\lambda_{0},\lambda_{1},\lambda_{2})} \xrightarrow{\mu} \tn{Gr}_{Z \subset G}^{e,(\lambda_{0},\lambda_{2})}.
\end{equation*}

Recall the classical convolution affine Grassmannian $\tn{Gr}_{G/Z,\Spd(C)} \tilde{\times} \tn{Gr}_{G/Z,\Spd(C)}$ for $G/Z$, which is equipped with a canonical isomorphism
\begin{equation}\label{eq:convGriso}
    \tn{Gr}_{G/Z,\Spd(C)} \tilde{\times} \tn{Gr}_{G/Z,\Spd(C)} \xrightarrow{\sim}  \tn{Gr}_{G/Z,\Spd(C)} \times \tn{Gr}_{G/Z,\Spd(C)}
\end{equation}
sending $(\cP_{0}, \cP_{1}, (f_{0},f_{1}))$ to $((\cP_{0}, f_{1} \circ f_{0}),(\cP_{1},f_{1})).$

There is a locally constant map
\begin{equation*}
    \tn{Gr}_{G/Z,\Spd(C)} \tilde{\times} \tn{Gr}_{G/Z,\Spd(C)} \to \Hck_{G/Z,\Spd(C)}^{\loc,2} \to \underline{\Hom}_{X}(u,Z)^{2},
 \end{equation*}
and for $\lambda, \lambda' \in \Hom_{F}(u,Z)$ we define $\tn{Gr}_{G/Z,\Spd(C)}^{(\lambda,\lambda')}$ as the preimage of $(\lambda,\lambda')$ under this map. Using this convention, we observe that convolution gives a map
\begin{equation*}
\tn{Gr}_{G/Z,\Spd(C)}^{(\lambda,\lambda')} \to \tn{Gr}_{G/Z,\Spd(C)}^{\lambda \lambda'}.
\end{equation*}

\begin{rque} The above displayed map is the same as applying the identification 
\begin{equation*}\tn{Gr}_{G/Z,\Spd(C)} \tilde{\times} \tn{Gr}_{G/Z,\Spd(C)} \xrightarrow{\eqref{eq:convGriso}} \tn{Gr}_{G/Z,\Spd(C)} \times \tn{Gr}_{G/Z,\Spd(C)}
\end{equation*}
and then, denoting the map $\tn{Gr}_{G/Z,\Spd(C)} \to \underline{\Hom}_{X}(u,Z)$ by $q$, taking the map $(q_{1}/q_{2}, q_{2})$ (where $q_{i}$ is $q$ composed with the $i$-projection).
\end{rque}


\begin{lem}\label{lem:Grcompconv} For fixed $\lambda_{0}, \lambda_{1},\lambda_{2} \in \Hom_{F}(u,Z)$, there is an $\mc{A}_{\lambda_{2}}$-equivariant equivalence 
\begin{equation*}
 \pi^{(\lambda_{0},\lambda_{1},\lambda_{2})}\colon \tn{Gr}_{Z \subset G}^{e,(\lambda_{0},\lambda_{1},\lambda_{2})} \to \tn{Gr}_{G/Z,\Spd(C)}^{(\lambda_{0}/\lambda_{1},\lambda_{1}/\lambda_{2})}
\end{equation*}
making the following diagram commute
    \begin{equation}\label{eq:Grconvdiag}
    \begin{tikzcd}
        \tn{Gr}_{Z \subset G}^{e,(\lambda_{0},\lambda_{1},\lambda_{2})} \arrow{r} \arrow["\pi^{(\lambda_{0},\lambda_{1},\lambda_{2})}"]{d} &   \tn{Gr}_{Z \subset G}^{e,(\lambda_{0},\lambda_{2})} \arrow{d} \\
       \tn{Gr}_{G/Z,\Spd(C)}^{(\lambda_{0}/\lambda_{1},\lambda_{1}/\lambda_{2})} \arrow{r} & \tn{Gr}_{G/Z,\Spd(C)}^{\lambda_{0}/
\lambda_{2}},
    \end{tikzcd}
    \end{equation}
    where the horizontal maps are convolution and the right-hand vertical map is from Proposition \ref{prop:mainGrisom}.
    \end{lem}

    \begin{proof}
We first define $\pi^{(\lambda_{0},\lambda_{1},\lambda_{2})}$ and prove it is an equivalence. It is given by sending the object $(\cP_{0},\cP_{1}, (f_{0}, f_{1})) \in \tn{Gr}_{Z \subset G}^{e,(\lambda_{0},\lambda_{1},\lambda_{2})} (S)$ to 
\begin{equation*}
(\ov{\cP_{0}},\ov{\cP_{1}},\bar{f}_{0} \circ \varphi_{\lambda_{1}/\lambda_{0},\cP_{0}},\ov{s(\lambda_{2})}_{*}(\tilde{\psi})\circ (\bar{f}_{1} \circ \varphi_{\lambda_{2}/\lambda_{1},\cP_{1}}))
\end{equation*}
where again we are using a fixed auxiliary $\bbD_{C}$-trivialization $\tilde{\psi}$ of $\mc{T}_{C}^{\loc}$. As is the case (the isomorphism \eqref{eq:convGriso}) for the classical convolution affine Grassmannian, combining convolution with the projection gives a canonical isomorphism
   \begin{equation}\label{eq:econvGriso}
        \tn{Gr}_{Z \subset G}^{e,(\lambda_{0},\lambda_{1},\lambda_{2})} \xrightarrow{\sim}  \tn{Gr}_{Z \subset G}^{e,(\lambda_{0},\lambda_{2})} \times \tn{Gr}_{Z \subset G}^{e,(\lambda_{1},\lambda_{2})};
   \end{equation}
   the proof of this isomorphism claim is straightforward and left to the reader.

   The fact that $\pi^{(\lambda_{0},\lambda_{1},\lambda_{2})}$ is an equivalence follows from the commutativity of the following diagram:
   \[
\begin{tikzcd}
     \tn{Gr}_{Z \subset G}^{e,(\lambda_{0},\lambda_{1},\lambda_{2})} \arrow["\eqref{eq:econvGriso}"]{r} \arrow["\pi^{(\lambda_{0},\lambda_{1},\lambda_{2})}"]{d} & \tn{Gr}_{Z \subset G}^{e,(\lambda_{0},\lambda_{2})} \times \tn{Gr}_{Z \subset G}^{e,(\lambda_{1},\lambda_{2})} \arrow{d} \\     \tn{Gr}_{G/Z,\Spd(C)}^{(\lambda_{0}/\lambda_{1},\lambda_{1}/\lambda_{2})} \arrow["\eqref{eq:convGriso}"]{r} & \tn{Gr}_{G/Z,\Spd(C)}^{\lambda_{0}/\lambda_{2}} \times \tn{Gr}_{G/Z,\Spd(C)}^{\lambda_{1}/\lambda_{2}},
\end{tikzcd}
   \]
   where all other maps are isomorphisms by the above arguments and Proposition \ref{prop:mainGrisom}.
   
   It remains to show that the diagram \eqref{eq:Grconvdiag} commutes. Going right and then down gives
   \begin{equation*}
       (\ov{\cP_{0}}, \ov{s(\lambda_{2})}_{*}(\tilde{\psi}) \circ \ov{z(\lambda_{1}/\lambda_{0},\lambda_{2}/\lambda_{1})_{*}(\psi^{-1}) \circ f_{1} \circ (\Q_{\lambda_{2}/\lambda_{1},S} \cdot f_{0})} \circ \varphi_{\lambda_{2}/\lambda_{0},\cP_{0}}).
   \end{equation*}
   Going the other direction yields
   \begin{equation*}
        (\ov{\cP_{0}},  \ov{s(\lambda_{2})}_{*}(\tilde{\psi}) \circ \bar{f}_{1} \circ \varphi_{\lambda_{2}/\lambda_{1},\cP_{1}} \circ \bar{f}_{0} \circ \varphi_{\lambda_{1}/\lambda_{0},\cP_{0}}),
   \end{equation*}
    and so again we are done by the definition of the cocycle $z$ in the same way as in the proof of Proposition \ref{prop:Hckconv} (the Hecke stack analogue of this statement).        
    \end{proof}


    \begin{rque}
        One can define $\tn{Gr}_{Z \subset G}^{e,(\lambda_{0}, \dots, \lambda_{n})}$ in a manner completely analogous to Definition \ref{defi:convGre}. The same proof as that of Lemma \ref{lem:Grcompconv} shows that there is an equivalence 
        \begin{equation*}
        \tn{Gr}_{Z \subset G}^{e,(\lambda_{0}, \dots, \lambda_{n})} \to \tn{Gr}_{G/Z,\Spd(C)}^{(\lambda_{0}/\lambda_{1}, \dots, \lambda_{n-1}/\lambda_{n})}
       \end{equation*}
        for any $n \geq 1$. Moreover, we have an identification 
        \begin{equation*}
      \mc{A}_{\lambda_{n}} \backslash \tn{Gr}_{Z \subset G}^{e,(\lambda_{0}, \dots, \lambda_{n})}  = \Hck_{Z \subset G}^{e,\loc,n,(\lambda_{0}, \dots, \lambda_{n})} 
        \end{equation*}
        which is compatible with any of the convolution maps for $1 \leq i \leq n-1$.
    \end{rque}
   

 For any small v-stack $S \to \Spd(C)$ we set $\Hck_{Z \subset G,S}^{e,\loc,1} := \Hck_{Z \subset G}^{e,\loc,1} \times_{\Spd(C)} S$, similarly with $\Hck_{Z \subset G,S}^{e,\loc,1,(\lambda_{0},\lambda_{1})}$ and $\tn{Gr}_{Z \subset G,S}^{e}$. One can define perverse sheaves on $\Hck_{Z \subset G,S}^{e,\loc,1}$ in a manner which is a verbatim analogue of \cite[Definition/Proposition VI.7.1]{Geometrization}, whose proof is identical because of the isomorphism of affine Grassmannians provided by Proposition \ref{prop:mainGrisom}:

 \begin{prop}
     Let $S \to \Spd(C)$ be a small v-stack. For any $\lambda_{0},\lambda_{1} \in \Hom_{F}(u,Z)$, there is a unique $t$-structure $({^{p}}D^{\leq0},{^{p}D}^{\geq 0})$ on $D_{\tn{\'{e}t}}(\Hck_{Z \subset G,S}^{e,\loc,1,(\lambda_{0},\lambda_{1})},\Lambda)$ such that 
     \begin{equation*}
     A \in {^{p}D}^{\leq 0}_{\tn{\'{e}t}}(\Hck_{Z \subset G,S}^{e,\loc,1,(\lambda_{0},\lambda_{1})},\Lambda)
     \end{equation*}
     if and only if for all geometric points $\bar{s}$ of $S$ and all $\nu \in X_{*}(T/Z)^{+,\lambda_{0}/\lambda_{1}}$ the pullback of $A$ to each open Schubert cell $\Hck_{Z \subset G,\bar{s},\nu}^{e,\loc,1,(\lambda_{0},\lambda_{1})}$ of $\Hck_{Z \subset G,\bar{s}}^{e,\loc,1,(\lambda_{0},\lambda_{1})}$ (as in Remark \ref{rque:strat}) sits in cohomological degrees $\leq - \langle 2\rho, \nu \rangle$, where $\rho$ is the half-sum of the positive roots of $G$.
 \end{prop}

 The above proposition, along with the decomposition 
 \begin{equation*}
 \Hck_{Z \subset G,S}^{e,\loc,1} = \bigsqcup_{(\lambda_{0},\lambda_{1}) \in \Hom_{F}(u,Z)^{2}}\Hck_{Z \subset G,S}^{e,\loc,1,(\lambda_{0},\lambda_{1})}
 \end{equation*}
 lets us define (relative) perverse sheaves on  $\Hck_{Z \subset G,S}^{e,\loc,1}$ for $S$ as above.

The following direct analogue of \cite[Proposition VI.7.2]{Geometrization} follows immediately:
 \begin{prop}\label{prop:FSVI.7.2}
 For any $S \to \Spd(C)$ a small v-stack the pullback functors 
 \begin{equation*}
 \tn{Perv}(\Hck_{Z \subset G,S}^{e,\loc,1,(\lambda_{0},\lambda_{1})},\Lambda) \to D_{\tn{\'{e}t}}(\tn{Gr}_{Z \subset G,S}^{e,(\lambda_{0},\lambda_{1})},\Lambda)^{\tn{bd}}
\end{equation*}
and
\begin{equation*}
 \tn{Perv}(\Hck_{Z \subset G,S}^{e,\loc,1},\Lambda) \to D_{\tn{\'{e}t}}(\tn{Gr}_{Z \subset G,S}^{e},\Lambda)^{\tn{bd}}
\end{equation*}
are fully faithful.
 \end{prop}

     \begin{corol}\label{prop:perveequiv1}
For a fixed $\lambda_{0}, \lambda_{1} \in \Hom_{F}(u,Z)$, the pullback along
\begin{equation*}
\Hck^{e,\loc,1,(\lambda_{0},\lambda_{1})}_{Z \subset G} \xrightarrow{\pi} \Hck^{\loc,1,\lambda_{0}/
\lambda_{1}}_{G/Z,\Spd(C)}
\end{equation*}
induces an equivalence of categories
\begin{equation}\label{eq:pervequivbis}
    \tn{Perv}(\Hck^{\loc,1,\lambda_{0}/
\lambda_{1}}_{G/Z,\Spd(C)},\Lambda) \to \tn{Perv}(\Hck^{e,\loc,1,(\lambda_{0},\lambda_{1})}_{Z \subset G},\Lambda).
\end{equation}
\end{corol}

\begin{proof}
   Proposition \ref{prop:mainGrisom} implies (cf. Corollary \ref{rque:sheafequiv}) that pullback along $\pi$ gives an equivalence 
    \begin{equation*}
    \tn{Perv}(\Hck^{\loc,1,\lambda_{0}/\lambda_{1}}_{G/Z,\Spd(C)},\Lambda) \xrightarrow{\sim} \tn{Perv}(\mc{A}_{\lambda_{1}} \backslash \tn{Gr}_{G/Z,\Spd(C)}^{\lambda_{0}/\lambda_{1}},\Lambda).
\end{equation*}

We claim that the pullback map
\begin{equation*}
  \tn{Perv}(L^{+}(G/Z) \backslash \tn{Gr}_{G/Z,\Spd(C)}^{\lambda_{0}/\lambda_{1}},\Lambda)  \to \tn{Perv}(\mc{A}_{\lambda_{1}} \backslash \tn{Gr}_{G/Z,\Spd(C)}^{\lambda_{0}/\lambda_{1}},\Lambda)
\end{equation*}
is also an equivalence; it suffices by Remark \ref{rque:Alambda} to prove this after identifying $\mc{A}_{\lambda_{1}}$ with $L^{+}G$. 

Having made this identification, this second asserted equivalence holds because, by Lemma \ref{lem:Grgerbe}, the $L^{+}G$-orbits on $\tn{Gr}_{G/Z,\Spd(C)}^{\lambda_{0}/\lambda_{1}}$
are the same as the $L^{+}(G/Z)$ orbits and for all points $X$ of $\tn{Gr}_{G/Z}(C)$ the subgroup $Z=L^{+}Z$ is contained in $(L^{+}G)_{X}^{\circ}$. To see this final claimed containment, note that $(L^{+}G)_{X}$ is the preimage of $(L^{+}(G/Z))_{X}$ in $L^{+}G$; we can reduce to the case where $X = [t^{\bar{\nu}}]$ with $t^{\bar{\nu}} \in L(T/Z)$ and $\bar{\nu} \in X_{*}(T/Z)$, so that $(L^{+}G)^{\circ}_{X}$ contains $L^{+}(T)$ and thus also $Z = L^{+}Z$. 
\end{proof}

\begin{defi}\label{defi:satake-category}
   We define the Satake category
    \begin{equation*}
    \tn{Sat}(\Hck_{Z \subset G}^{e,\loc,1,(\lambda_{0},\lambda_{1})},\Lambda) \subset D_{\tn{\'{e}t}}(\tn{Gr}_{Z \subset G}^{e,(\lambda_{0},\lambda_{1})},\Lambda)
    \end{equation*}
    to be the preimage of $\tn{Sat}(\Hck_{G/Z,\Spd(C)}^{\loc,1,\lambda_{0}/\lambda_{1}},\Lambda)$ (in the sense of \cite[Definition VI.7.8]{Geometrization})
    under pullback by the isomorphism $\pi^{(\lambda_{0},\lambda_{1})}$. 
    
    \begin{rque} It is straightforward to check (using Proposition \ref{prop:mainGrisom}) that this is the same as the full subcategory of all objects in $\mc{D}_{\tn{\'{e}t}}(\Hck_{Z \subset G}^{e,\loc,1,(\lambda_{0},\lambda_{1})},\Lambda)^{\tn{bd}}$ that are universally locally acyclic and flat perverse. 
    \end{rque}
    
    Similarly, we define
    \begin{equation*}
\tn{Sat}(\Hck_{Z \subset G}^{e,\loc,1},\Lambda) = \prod_{(\lambda_{0},\lambda_{1}) \in \Hom_{F}(u,Z)^{2}} \tn{Sat}(\Hck_{Z \subset G}^{e,\loc,1,(\lambda_{0},\lambda_{1})},\Lambda)  \subset \mc{D}_{\tn{\'{e}t}}(\Hck_{Z \subset G}^{e,\loc,1},\Lambda)^{\tn{bd}},
    \end{equation*}
    which again we can characterize as all bounded objects that are universally locally acyclic and flat perverse. 
\end{defi}

 \begin{prop}\label{prop:pervconv} Fix $\lambda_{0}, \lambda_{1},\lambda_{2} \in \Hom_{F}(u,Z)$.
     \begin{enumerate}
         \item{If $A \in {^{p}D}^{\leq 0}_{\tn{\'{e}t}}(\Hck_{Z \subset G}^{e,\loc,1,(\lambda_{0},\lambda_{1})},\Lambda)^{\tn{bd}}$ and $A' \in {^{p}D}^{\leq 0}_{\tn{\'{e}t}}(\Hck_{Z \subset G}^{e,\loc,1,(\lambda_{1},\lambda_{2})},\Lambda)^{\tn{bd}}$ then $A \star A' \in {^{p}D}^{\leq 0}_{\tn{\'{e}t}}(\Hck_{Z \subset G}^{e,\loc,1,(\lambda_{0},\lambda_{2})},\Lambda)^{\tn{bd}}$.}
\item{If $A,A' \in \tn{Sat}(\Hck_{Z \subset G}^{e,\loc,1},\Lambda)$ then $A \star A' \in \tn{Sat}(\Hck_{Z \subset G}^{e,\loc,1},\Lambda)$.}
         \item{There is a commutative convolution diagram
    \[
    \begin{tikzcd}
       \tn{Sat}(\Hck^{e,\loc,1,(\lambda_{0},\lambda_{1})}_{Z \subset G},\Lambda)\times \tn{Sat}( \Hck^{e,\loc,1,(\lambda_{1},\lambda_{2})}_{Z \subset G},\Lambda) \arrow{r} \arrow["\sim"]{d} &  \tn{Sat}(\tn{Hck}_{Z \subset G}^{e,\loc,1,(\lambda_{0},\lambda_{2})},\Lambda) \arrow["\sim"]{d} \\
        \tn{Sat}(\Hck^{\loc,1,\lambda_{0}/
\lambda_{1}}_{G/Z,\Spd(C)},\Lambda) \times \tn{Sat}(\Hck^{\loc,1,\lambda_{1}/
\lambda_{2}}_{G/Z,\Spd(C)},\Lambda) \arrow{r}  & \tn{Sat}(\Hck^{\loc,1,\lambda_{0}/\lambda_{2}
}_{G/Z,\Spd(C)},\Lambda).
\end{tikzcd}
\]}

     \end{enumerate}
 \end{prop}

 \begin{proof}
     Combining Lemma \ref{lem:outer-faces} with Corollary \ref{prop:perveequiv1} and then using the commutative diagram
     \[
     \begin{tikzcd}
     \Hck_{Z \subset G}^{e,\loc,2,(\lambda_{0},\lambda_{1},\lambda_{2})} \arrow{d} \arrow{r} & \Hck_{Z \subset G}^{e,\loc,1,(\lambda_{0},\lambda_{2})} \arrow{d} \\
\Hck_{G/Z,\Spd(C)}^{\loc,2,(\lambda_{0}/\lambda_{1},\lambda_{1}/\lambda_{2})} \arrow{r} & \Hck_{G/Z,\Spd(C)}^{\loc,1,\lambda_{0}/\lambda_{2}},
     \end{tikzcd}
     \]
     which is Cartesian by Proposition \ref{prop:mainGrisom} and Lemma \ref{lem:Grcompconv}, implies that convolution of perverse sheaves on $\Hck_{Z \subset G}^{e,1,\loc}$ can be computed using the analogous convolution on $\Hck_{G/Z,\Spd(C)}^{\loc,1}$, giving statement (3).  All of the remaining claims therefore follow from \cite[Proposition VI.8.1]{Geometrization}.
 \end{proof}


Passing to the dual side, denote by $\widehat{Z}$ the $\Lambda$-group scheme Cartier dual to $\Hom_{\Z}(X^{*}(Z), \mathbb{Q}/\Z)$, which fits into the central extension
\begin{equation*}
    1 \to \widehat{Z} \to \widehat{G/Z} \to \widehat{G} \to 1.
\end{equation*}
 Note that there is an identification
\begin{equation*}
    X^{*}(\widehat{Z}) \xrightarrow{\sim} \Hom_{\ov{F}}(\mu,Z) \xrightarrow{\sim} \Hom_{F}(u,Z),
\end{equation*}
which we will use without comment. 

By taking the weight decomposition of $V \in \Rep_{\Lambda}(\widehat{G/Z})$ and restricting to $\widehat{Z}$ (viewed as a $\Lambda$-group scheme), this category has a natural $X^{*}(\widehat{Z})$-grading.

Consider the category 
\begin{equation*}
\bigoplus_{(\lambda_{0}, \lambda_{1}) \in \Hom_{F}(u, Z)^{2}} \tn{Rep}_{\Lambda}(\widehat{G/Z})^{\lambda_{0}/\lambda_{1}},
\end{equation*}
which carries an additive, associative, non-symmetric ``convolution'' monoidal structure determined by the formula (defined on individual summands)
\begin{equation*}
v_{(\lambda_{i},\lambda_{j})} \star w_{(\lambda_{k}, \lambda_{l})} = (v \otimes w)_{(\lambda_{i},\lambda_{l})}
\end{equation*}
if $\lambda_{j} =\lambda_{k}$ and zero otherwise.

\begin{rque}\label{rque:altmonoidal}
One can identify the monoidal category $\bigoplus_{(\lambda_{0}, \lambda_{1}) \in \Hom_{F}(u, Z)^{2}} \tn{Rep}_{\Lambda}(\widehat{G/Z})^{\lambda_{0}/\lambda_{1}}$ (with $\star$) with the category 
\begin{equation*}
\Coh^{\heartsuit, \Lambda-\tn{flat}}(\pt/\widehat{G/Z} \times_{\pt/\widehat{G}} \pt/\widehat{G/Z})
\end{equation*}
equipped with the convolution monoidal structure. This gives an alternative interpretation of the monoidal category appearing above.
\end{rque}


We obtain an analogue of the geometric Satake isomorphism:

\begin{thm}\label{thm:satake}
    We have an equivalence of monoidal categories
    \begin{equation}\label{eq:pervequiv}
        \tn{Sat}(\Hck^{e,\loc,1}_{Z \subset G},\Lambda) \xrightarrow{\sim} (\bigoplus_{(\lambda_{0}, \lambda_{1}) \in \Hom_{F}(u, Z)^{2}} \tn{Rep}_{\Lambda}(\widehat{G/Z})^{\lambda_{0}/\lambda_{1}},\star).
    \end{equation}
\end{thm}

\begin{proof}
    Using the identification of monoidal categories
    \begin{equation*}
\tn{Sat}(\Hck^{\loc,1}_{G/Z,\Spd(C)},\Lambda) \xrightarrow{\sim} \tn{Rep}_{\Lambda}(\widehat{G/Z}),
    \end{equation*}
    from \cite[Theorem I.6.3]{Geometrization} we observe that the equivalence \eqref{eq:pervequivbis} from Corollary \ref{prop:perveequiv1} is identified with an equivalence of categories
     \begin{equation*}
\tn{Sat}(\Hck^{e,\loc,1,(\lambda_{0},
\lambda_{1})}_{Z \subset G},\Lambda) \xrightarrow{\sim} \tn{Rep}_{\Lambda}(\widehat{G/Z})^{\lambda_{0}/\lambda_{1}}.
    \end{equation*}
 
    This map is given explicitly by sending a representation $V \in \tn{Rep}_{\Lambda}(\widehat{G/Z})^{\lambda_{0}/\lambda_{1}}$ to the pullback of the Satake sheaf $\mathcal{S}_{V} \in \tn{Sat}(\Hck_{G/Z}^{\loc,1,\lambda_{0}/\lambda_{1}},\Lambda)$ to $\tn{Sat}(\Hck_{Z \subset G}^{e,\loc,1,(\lambda_{0},\lambda_{1})},\Lambda)$.

    Moreover, Proposition \ref{prop:pervconv} implies that, via the above identification, the convolution map
    \begin{equation*}
\tn{Sat}(\Hck^{e,\loc,1,(\lambda_{0},
\lambda_{1})}_{Z \subset G},\Lambda)  \times \tn{Sat}(\Hck^{e,\loc,1,(\lambda_{1},
\lambda_{2})}_{Z \subset G},\Lambda) \to \tn{Sat}(\Hck^{e,\loc,1,(\lambda_{0},
\lambda_{2})}_{Z \subset G},\Lambda) 
    \end{equation*}
    is identified with the tensor product map
    \begin{equation*}
        \tn{Rep}_{\Lambda}(\widehat{G/Z})^{\lambda_{0}/\lambda_{1}} \times \tn{Rep}_{\Lambda}(\widehat{G/Z})^{\lambda_{1}/\lambda_{2}} \to \tn{Rep}_{\Lambda}(\widehat{G/Z})^{\lambda_{0}/\lambda_{2}}.
    \end{equation*}
    Now applying the disjoint union decomposition
    \begin{equation*}
\Hck^{e,\loc,1}_{Z \subset G} = \bigsqcup_{(\lambda_{0}, \lambda_{1}) \in \Hom_{F}(u, Z)^{2}} \Hck^{e,\loc,1,(\lambda_{0},\lambda_{1})}_{Z \subset G}
\end{equation*}
from above, we see that convolution between perverse sheaves supported on $(\lambda_{i},\lambda_{j})$ and $(\lambda_{k}, \lambda_{l})$-components with $\lambda_{j} \neq \lambda_{k}$ yields zero. Together with our previous observations from this proof, we deduce an identification
\begin{equation*}
     \tn{Sat}(\Hck^{e,\loc,1}_{Z \subset G},\Lambda) \xrightarrow{\sim} \bigoplus_{(\lambda_{0}, \lambda_{1}) \in \Hom_{F}(u, Z)^{2}} \tn{Rep}_{\Lambda}(\widehat{G/Z})^{\lambda_{0}/\lambda_{1}},
\end{equation*}
with the monoidal structure $\star$ on the right-hand side. 
\end{proof}

It is also clear that, via the above identification, the pullback map 
\begin{equation*}
    \tn{Sat}(\Hck_{G/Z,\Spd(C)}^{\loc,1},\Lambda) \to \tn{Sat}(\Hck^{e,\loc,1}_{Z \subset G},\Lambda)
\end{equation*}
corresponds to the embedding
\begin{equation}\label{eq:G/Zpullback}
\tn{Rep}_{\Lambda}(\widehat{G/Z}) \hookrightarrow \bigoplus_{(\lambda_{0}, \lambda_{1}) \in \Hom_{F}(u, Z)^{2}} \tn{Rep}_{\Lambda}(\widehat{G/Z})^{\lambda_{0}/\lambda_{1}}
\end{equation}
determined by sending $\pi \in \tn{Rep}_{\Lambda}(\widehat{G/Z})^{\lambda}$ to $\bigoplus_{(\lambda' \lambda,\lambda') \in \Hom_{F}(u,Z)^{2}} \pi$. 

\begin{rque}
Using the alternative description from Remark \ref{rque:altmonoidal}, the embedding \eqref{eq:G/Zpullback} corresponds to pushforward of coherent sheaves along the diagonal
\begin{equation*}
\pt/(\widehat{G/Z}) \xrightarrow{\Delta} \pt/(\widehat{G/Z}) \times_{\pt/\widehat{G}} \pt/(\widehat{G/Z}).
\end{equation*}
\end{rque}

\begin{lem}
    The embedding \eqref{eq:G/Zpullback} is a monoidal functor.
\end{lem}

\begin{proof}
    This is an elementary calculation.
\end{proof}

The following result comes from the proof of Theorem \ref{thm:satake} and the functoriality of the geometric Satake isomorphism from \cite{Geometrization}:

\begin{corol}\label{corol:satZfunc} The isomorphism from Theorem \ref{thm:satake} is functorial in $Z \subset Z' \subset Z_{G}$ via pullback along the transition maps \eqref{eq:Ztrans} and the functor
\begin{equation*}
    \bigoplus_{(\lambda_{0}, \lambda_{1}) \in \Hom_{F}(u, Z)^{2}} \tn{Rep}_{\Lambda}(\widehat{G/Z})^{\lambda_{0}/\lambda_{1}}\to \bigoplus_{(\lambda_{0}, \lambda_{1}) \in \Hom_{F}(u, Z')^{2}} \tn{Rep}_{\Lambda}(\widehat{G/Z'})^{\lambda_{0}/\lambda_{1}}
\end{equation*}
induced by the maps $\widehat{G/Z'} \to \widehat{G/Z}$ and $Z \hookrightarrow Z'$.
\end{corol}

\subsection{Global Hecke stacks}

Recall from \cite[Proposition 19.1.2]{SW20} that, for any $S \in \Perfd_{C}$ and effective Cartier divisor $y' \in \Div^{d}(S)$, there is a faithful functor from the category of triples $(U,V,f)$ consisting of a $G_{S}$-torsor $U$ on $X_{S} \setminus \Gamma_{y'}$, a $G_{S}$-torsor $V$ on $\bbD_{y'}$, and an isomorphism 
\begin{equation*}
f \colon U|_{\bbD_{y'}^{\circ}} \xrightarrow{\sim} V|_{\bbD_{y'}^{\circ}}
\end{equation*}
to the category of $G_{S}$-torsors on $X_{S}$ (see also \cite[\S III.3]{Geometrization}). We call such a triple $(U,V,f)$ a \textbf{Beauville--Laszlo triple} for $X_{S}$. One can define the category of Beauville-Laszlo triples in a completely analogous manner for $[\mc{T}_{S}/\tilde{t}_{S}]$ as well.

\begin{lem}\label{lem:GerbeBL}
There is a functor from the category of Beauville-Laszlo triples for $[\mc{T}_{S}/\tilde{t}_{S}]$ to the category of $G_{S}$-torsors on $[\mc{T}_{S}/\tilde{t}_{S}]$.
\end{lem}

\begin{proof}
By Remark \ref{rque:fingerbe}.(1) we can replace $[\mc{T}_{S}/\tilde{t}_{S}]$ and $[\mc{T}_{S}^{\loc}/\tilde{t}_{S}]$ with $\mc{E}:=[\mc{T}_{E/F,S}/\tilde{t}_{E/F,n,S}]$ and $\mc{E}^{\loc} := [\mc{T}_{E/F,S}^{\loc}/\tilde{t}_{E/F,n,S}]$ for some finite Galois $E/F$ and $n \in \mathbb{N}$. By Remark \ref{rque:fingerbe}.(2) the gerbe $\mc{E}$ is trivialized by the finite \'{e}tale cover $X_{S}^{(E')}$ for some finite Galois extension $E'/F$ and is represented by the \v{C}ech $2$-cocycle $\xi_{\mc{E}} := \xi_{E/F,n}$, which are compatible for varying $E/F$ and $n$.

According to \cite[Proposition 2.50]{Dillery23}, the category of $G_{S}$-torsors on $\mc{E}$ is equivalent to the category of \textbf{$\xi_{\mc{E}}$-twisted $G_{S}$-torsors} whose objects are pairs $(T,\phi)$, where $T$ is a $u_{E/F,n}$-equivariant $G_{S}$-torsor on $X_{S}^{(E')}$ and $\pi_{1}^{*}T \xrightarrow{\phi} \pi_{2}^{*}T$ is an equivariant isomorphism of $G_{S}$-torsors on $X_{S}^{(E')} \times_{X_{S}} X_{S}^{(E')}$ such that $d\phi$ (as defined in \cite[Lemma 2.37]{Dillery23}) is translation by $\xi_{\mc{E}}$. Morphisms are isomorphisms of torsors that are compatible with the respective maps $\phi$.

Given a Beauville-Laszlo triple $(\mc{U},\mc{V},f)$ for $\mc{E}$ (descended from a given triple for $[\mc{T}_{S}/\tilde{t}_{S}]$), we denote by $((U,\phi_{U}),(V,\phi_{V}),g)$ the corresponding triple where $(U,\phi_{U})$ and $(V,\phi_{V})$ are twisted torsors.

The triple $(U,V,g)$ is a Beauville-Laszlo triple for $X_{S}^{(E')}$, so by Beauville-Laszlo for $X_{S}^{(E')}$ we can glue this to a $G_{S}$-torsor $T$ over $X_{S}^{(E')}$. Moreover, the functoriality of Beauville-Laszlo for $X_{S}^{(E')}$, after applying the identification 
\begin{equation*}
X_{S}^{(E')} \times_{X_{S}} X_{S}^{(E')} = \bigsqcup_{\Gal_{E'/F}} X_{S}^{(E')},
\end{equation*}
implies that, since $\phi_{V} = \phi_{U} \circ \pi_{1}^{*}g$ on $\bigsqcup\bbD_{S}^{(E'),\circ}$, we can glue these two maps to obtain an isomorphism of torsors 
\begin{equation*}
    \pi_{1}^{*}T \xrightarrow{\phi} \pi_{2}^{*}T
\end{equation*}
over $X_{S}^{(E')} \times_{X_{S}} X_{S}^{(E')}$. Finally, faithfulness of Beauville-Laszlo for $X_{S}^{(E')}$ implies that since $d\phi_{V}$ and $d\phi_{U}$ are translation by $\xi_{\mc{E}}$, so is $d\phi$. The pair $(T,\phi)$ gives the desired $G_{S}$-torsor on $\mc{E}$. All of these constructions are functorial in the triple $(\mc{U},\mc{V},f)$ because of the functoriality of Beauville-Laszlo for $X_{S}^{(E')}$. 
\end{proof}

Consider the Hecke correspondence
\[
\begin{tikzcd}
  &  \Hck_{Z \subset G}^{e,1} \arrow["\overleftarrow{h}"]{ld} \arrow["\overrightarrow{h}"]{rd} & \\
  \Bun_{G}^{e} & & \Bun_{G}^{e}.
\end{tikzcd}
\]

\begin{prop}\label{prop:Heckeproper} The maps $\overleftarrow{h}$ and $\overrightarrow{h}$ for $\Hck^{e,1}_{Z \subset G}$ are ind-proper.
\end{prop}
\begin{proof} It suffices to prove this for each $\Hck^{e,1,(\lambda_{0},\lambda_{1})}_{Z \subset G}$. In this setting, pulling back the closed Schubert cells of $\Hck_{Z \subset G}^{e,\loc,1}$ (cf. Remark \ref{rque:strat}) gives (as a v-stack)
\begin{equation*}
\Hck^{e,1,(\lambda_{0},\lambda_{1})}_{Z \subset G} = \varinjlim_{\nu} \Hck^{e,1,(\lambda_{0},\lambda_{1})}_{Z \subset G,\leq \nu} 
\end{equation*}
for $\nu \in X_{*}(T/Z)^{+,\lambda_{0}/\lambda_{1}}$. 

Moreover, it suffices to show that $\overrightarrow{h}$ is proper for each $\Hck^{e,1,(\lambda_{0},\lambda_{1})}_{Z \subset G,\leq \nu}$, since $\overleftarrow{h}$ can be identified with $\overrightarrow{h}$ (for $\Hck_{Z \subset G}^{e,1,(\lambda_{1},\lambda_{0})}$) after applying the isomorphism given by switching the torsors.

Working v-locally on the target allows us to check the result for $S \to \Bun_{G}^{e,\lambda_{1}}$ strictly totally disconnected. We claim that the fibers of $\overrightarrow{h}$ over points $\mc{X} \in \Bun_{G}^{e,\lambda_{1}}(S)$ for such $S \in \Perfd_{C}$ are isomorphic to the extended twisted affine Grassmannian
\begin{equation*}
\tn{Gr}_{Z \subset G,\leq \nu}^{e,(\lambda_{0},\lambda_{1})} \tilde{\times} \mc{X}
\end{equation*}
parametrizing, for $S' \to S$, pairs $(\cP,f)$ of a $G_{S'}$-torsor $\cP$ on $[\mc{T}_{S'}^{\loc}/\tilde{t}_{S'}]$ with $\lambda_{\cP} = \lambda_{0}$ and $f$ an isomorphism from $\Q_{\lambda_{1}/\lambda_{0},S'} \cdot \cP$ to $\mc{X}$ on $[\mc{T}_{S'}^{\loc,\circ}/\tilde{t}_{S'}]$ bounded by $\nu$.
We could then deduce the result by Proposition \ref{prop:mainGrisom}, since identifying $\mc{X}$ with $\widetilde{\Q}_{\lambda_{1},S}$ over $[\mc{T}_{S}^{\loc}/\tilde{t}_{S}]$ (possible by \cite[Lemma III.2.6]{Geometrization} because $S$ is strictly totally disconnected) identifies $\tn{Gr}_{Z \subset G,\leq \nu}^{e,(\lambda_{0},\lambda_{1})} \tilde{\times} \mc{X}$ with $\tn{Gr}_{Z \subset G,\leq \nu}^{e,(\lambda_{0},\lambda_{1})}\times S$. 

To prove this description of the fiber, the same argument as for $\Hck_{G}$ holds provided that one has a surjective ``Beauville-Laszlo'' morphism
\begin{equation*}
    \tn{Gr}_{Z \subset G}^{e,(\lambda_{0},\lambda_{1})} \tilde{\times} \mc{X} \to \Bun_{G}^{e,\lambda_{0}} \times S,
\end{equation*}
of v-stacks. 

The existence of the morphism comes from Lemma \ref{lem:GerbeBL} because a given object $(\cP,f) \in [\tn{Gr}_{Z \subset G}^{e,(\lambda_{0},\lambda_{1})} \tilde{\times} \mc{X}](S)$ defines a Beauville-Laszlo triple
\begin{equation*}
(\cP, \Q_{\lambda_{0}/\lambda_{1},S}^{-1} \cdot\mc{X},\Q_{\lambda_{0}/\lambda_{1},S}^{-1}\cdot f)
\end{equation*}
for $[\mc{T}_{S}/\tilde{t}_{S}]$.
This morphism is surjective by the argument in \cite[Proposition III.3.1]{Geometrization}.
\end{proof}

\subsection{Extended Hecke action}\label{sec:Extended-Hecke-Action}
This subsection is mostly (until the main theorems at the end) a formal summary of some of the ideas in \cite[\S IX]{Geometrization} and why they still work in this extended setting. In this section, $\Lambda$ is a $\ov{\Z_{\ell}}$-algebra.

First, taking the inverse limit of the isomorphism from Theorem \ref{thm:satake} over all $\Z/\ell^{r}\Z[\sqrt{q}]$ gives an exact $\Z/\ell\Z[\sqrt{q}]$-linear monoidal functor
\begin{equation*}
    \bigoplus_{(\lambda_{0}, \lambda_{1}) \in \Hom_{F}(u, Z)^{2}} \tn{Rep}_{\Z_{\ell}[\sqrt{q}]}(\widehat{G/Z})^{\lambda_{0}/\lambda_{1}} \to \tn{Sat}(\Hck_{Z \subset G}^{e,\loc,1}, \Z_{\ell}[\sqrt{q}]),
\end{equation*}
where the right-hand side is the inverse limit of $\tn{Sat}(\Hck_{Z \subset G}^{e,\loc,1}, \Z/\ell^{r}\Z[\sqrt{q}])$ for all $r$. We pass everything further to $\ov{\Z_{\ell}}$ for simplicity.

Denote by $D_{\blacksquare}(\Bun_{G}^{e},\Lambda)$ the category of solid sheaves of $\Lambda$-modules on the v-site of $\Bun_{G}^{e}$. Recall from \cite[\S IX.2]{Geometrization} that there is a $\tn{Rep}_{\Lambda}(Q)$-linear monoidal functor
\begin{equation*}
\tn{Rep}_{\ov{\Z_{\ell}}}(\widehat{G} \rtimes Q) \to D_{\blacksquare}(\Hck_{G}^{\loc,1}, \ov{\Z_{\ell}})
\end{equation*}
given by composing the functor $V \mapsto \mathcal{S}_{V}$ with $A \mapsto \mathbb{D}(A)^{\vee}$, where $\bbD(A)$ denotes the (relative) Verdier dual in the sense of \cite[\S V.6]{Geometrization}.

In our setting, we obtain an analogous monoidal functor
\begin{equation*}
\bigoplus_{(\lambda_{0}, \lambda_{1}) \in \Hom_{F}(u, Z)^{2}} \tn{Rep}_{\ov{\Z_{\ell}}}(\widehat{G/Z})^{\lambda_{0}/\lambda_{1}} \to \tn{Sat}(\Hck_{Z \subset G}^{e,\loc,1},\ov{\Z_{\ell}})
\end{equation*}
and then composing with $A \mapsto \mathbb{D}(A)^{\vee}$, where $\mathbb{D}$ is the Verdier dual relative to the projection $\Hck_{Z \subset G}^{e,\loc,1} \to \Spd(C)/L^{+}G$
gives an exact monoidal functor
\begin{equation}\label{eq:solidmon}
\bigoplus_{(\lambda_{0}, \lambda_{1}) \in \Hom_{F}(u, Z)^{2}} \tn{Rep}_{\ov{\Z_{\ell}}}(\widehat{G/Z})^{\lambda_{0}/\lambda_{1}}  \to D_{\blacksquare}(\Hck_{Z \subset G}^{e,\loc,1},\ov{\Z_{\ell}}),
\end{equation}
where convolution in $D_{\blacksquare}(\Hck_{Z \subset G}^{e,\loc,1},\ov{\Z_{\ell}})$ is as defined in \cite[\S VII.5]{Geometrization} (all of the machinery developed in this aforementioned section of \cite{Geometrization} works for $\Sat(\Hck_{G}^{\loc,1},\ov{\Z_{\ell}})$ replaced with $\Sat(\Hck_{Z \subset G}^{e,\loc,1,(\lambda_{0},\lambda_{1})},\ov{\Z_{\ell}})$ or $\Sat(\Hck_{Z \subset G}^{e,\loc,1},\ov{\Z_{\ell}})$). 

We can then extend linearly to obtain an exact $\Lambda$-linear monoidal functor
\begin{equation*}
    \bigoplus_{(\lambda_{0}, \lambda_{1}) \in \Hom_{F}(u, Z)^{2}} \tn{Rep}_{\Lambda}(\widehat{G/Z})^{\lambda_{0}/\lambda_{1}}  \to D_{\blacksquare}(\Hck_{Z \subset G}^{e,\loc,1},\Lambda);
\end{equation*}
we denote the image of $V^{e}$ under this functor by $\mc{S}_{V^{e}}'$.

\begin{rque}
    Since each $\Hck_{Z \subset G}^{e,\loc,1,(\lambda_{0},\lambda_{1})} \to \Hck_{G/Z,\Spd(C)}^{\loc,1,\lambda_{0}/\lambda_{1}}$ is a $Z$-gerbe and we do not exclude the case where $\ell \mid |Z|$, it is not automatically true that the above functor \eqref{eq:solidmon} (and therefore its extension) is the same as pulling back the corresponding functor for $\Hck_{G/Z}^{\loc,1}$ along this gerbe. This potential discrepancy is due to the failure of taking Verdier dual to commute with pullback. 
\end{rque}

We then obtain a monoidal functor
\begin{equation}\label{eq:monoidal1}
    \bigoplus_{(\lambda_{0}, \lambda_{1}) \in \Hom_{F}(u, Z)^{2}} \tn{Rep}_{\Lambda}(\widehat{G/Z})^{\lambda_{0}/\lambda_{1}}  \to D_{\blacksquare}(\Hck_{Z \subset G}^{e,1},\Lambda)
\end{equation}
by pullback along $\Hck_{Z \subset G}^{e,1} \xrightarrow{\varepsilon} \Hck_{Z \subset G}^{e,\loc,1}$ (using Proposition \ref{prop:hckloc} for the monoidality).

Observe that there are two natural ``Hecke correspondence'' diagrams. The first is
\[
\begin{tikzcd}
  &  \Hck_{Z \subset G}^{e,1,(\lambda_{0},\lambda_{1})} \arrow["\overleftarrow{h}"]{ld} \arrow["\overrightarrow{h}"]{rd} & \\
  \Bun_{G}^{e,\lambda_{0}} & & \Bun_{G}^{e,\lambda_{1}},
\end{tikzcd}
\]
which gives rise to a $\Lambda$-linear functor
\begin{equation*}
    D_{\blacksquare}(\Hck_{Z \subset G}^{e,1,(\lambda_{0},\lambda_{1})},\Lambda) \to \tn{Fun}_{\Lambda}(D_{\blacksquare}(\Bun_{G}^{e,\lambda_{0}} \times \Spd(C),\Lambda), D_{\blacksquare}(\Bun_{G}^{e,\lambda_{1}} \times \Spd(C),\Lambda)).
\end{equation*}
Combining this across all $\lambda_{0},\lambda_{1}$ gives an analogous monoidal functor 
\begin{equation}\label{eq:monoidal2}
D_{\blacksquare}(\Hck_{Z \subset G}^{e,1},\Lambda) \to \tn{End}_{\Lambda}(D_{\blacksquare}(\Bun_{G}^{e} \times \Spd(C),\Lambda)).
\end{equation}

We thus obtain the composite functor, for any $V^{e} \in   \bigoplus_{(\lambda_{0}, \lambda_{1}) \in \Hom_{F}(u, Z)^{2}} \tn{Rep}_{\Lambda}(\widehat{G/Z})^{\lambda_{0}/\lambda_{1}}$:
\begin{equation*}
T^{e}_{V^{e}} \colon \mc{D}_{\tn{lis}}(\Bun_{G}^{e},\Lambda) \to D_{\blacksquare}(\Bun_{G}^{e} \times \Spd(C), \Lambda)
\end{equation*}
given by the formula
\begin{equation*}
    A \mapsto \overrightarrow{h}_{\natural}(\overleftarrow{h}^{*}A \prescript{\mathsmaller{\blacksquare}}{}\otimes_{\Lambda}^{L} (\varepsilon^{*}\mathcal{S}_{V^{e}}')),
\end{equation*}
where $\overrightarrow{h}_{\natural}$ is the left adjoint to $\overrightarrow{h}^{*}$ in the category of solid complexes of $\Lambda$-modules.


The following result is an analogue of \cite[Proposition IX.2.1]{Geometrization}:
\begin{prop}
    For any $V^{e} \in  \bigoplus_{(\lambda_{0}, \lambda_{1}) \in \Hom_{F}(u, Z)^{2}} \tn{Rep}_{\Lambda}(\widehat{G/Z})^{\lambda_{0}/\lambda_{1}}$, the functor $T^{e}_{V^{e}}$ defined above restricts to a functor 
    \begin{equation*}
T^{e}_{V^{e}} \colon \mc{D}_{\tn{lis}}(\Bun_{G}^{e}, \Lambda) \to \mc{D}_{\tn{lis}}(\Bun_{G}^{e}, \Lambda).
    \end{equation*}
    The resulting functor
    \begin{equation*}
         \bigoplus_{(\lambda_{0}, \lambda_{1}) \in \Hom_{F}(u, Z)^{2}} \tn{Rep}_{\Lambda}(\widehat{G/Z})^{\lambda_{0}/\lambda_{1}}  \to \tn{End}_{\Lambda}(\mc{D}_{\tn{lis}}(\Bun_{G}^{e},\Lambda)), \hspace{1mm} V^{e} \mapsto T^{e}_{V^{e}},
    \end{equation*}
    is monoidal.
\end{prop}

\begin{proof}
    For the first part, the proof of \cite[Proposition IX.2.1]{Geometrization} holds verbatim if one replaces $\Hck_{G}^{\loc,1}$ (in our notation---the notation \emph{loc. cit.} is different) with $\Hck_{Z \subset G}^{e,\loc,1,(\lambda_{0},\lambda_{1})}$ and uses Proposition \ref{prop:Heckeproper} and the description of the fibers of $\overrightarrow{h}$ from its proof along with the ``extended'' analogue of Demazure resolutions (as in \cite[\S VI.5]{Geometrization}) for the closure of Schubert cells of $\tn{Gr}^{e,(\lambda_{0},\lambda_{1})}_{Z \subset G}$ obtained by pulling back the usual Demazure resolutions for cells (cf. Remark \ref{rque:strat}) of $\tn{Gr}^{\lambda_{0}/\lambda_{1}}_{G/Z,\Spd(C)}$ along the isomorphism $\pi^{(\lambda_{0},\lambda_{1})}$ from Proposition \ref{prop:mainGrisom}. We are also using Corollary \ref{cor:Dlise} for the equivalence between $\mc{D}_{\tn{lis}}(\Bun_{G}^{e},\Lambda)$ and $\mc{D}_{\tn{lis}}(\Bun_{G}^{e} \times \Spd(C),\Lambda)$.

    The second part follows from the fact that $T^{e}_{V^{e}}$ is  obtained by restricting the endomorphism of $D_{\blacksquare}(\Bun_{G}^{e} \times \Spd(C),\Lambda)$ obtained from mapping $V^{e}$ along the composition \eqref{eq:monoidal1} with \eqref{eq:monoidal2}, which is monoidal, to an endomorphism of $\mc{D}_{\tn{lis}}(\Bun_{G}^{e},\Lambda)$.
\end{proof}

The second Hecke correspondence comes from the ``naive'' extended Hecke stacks (Definition \ref{defi:hcknaive}) $\Hck_{Z \subset G}^{e,(I_{i}),\tn{naive}}$ and $\Hck_{Z \subset G}^{e,\loc,(I_{i}),\tn{naive}}$ for a finite set $I$. Explicitly, it is given by 
\[
\begin{tikzcd}
  &  \Hck_{Z \subset G}^{e,(I_{i}),\tn{naive}} \arrow{ld} \arrow{rd} & \\
  \Bun_{G}^{e} & & \Bun_{G}^{e} \times (\Div^{1})^{I}.
\end{tikzcd}
\]

Recall from Proposition \ref{prop:naiveSA} that there is, for each $V \in \tn{Rep}_{\Lambda}((\widehat{G} \rtimes Q)^{I})$, a functor
\begin{equation*}
T^{e}_{V} \colon  \mc{D}_{\tn{lis}}(\Bun_{G}^{e},\Lambda) \to \mc{D}_{\tn{lis}}(\Bun_{G}^{e},\Lambda)^{B\Weil_{F}^{I}}.
\end{equation*}
These functors are defined by running the \cite{Geometrization} machinery for  a completely identical recipe as in \cite{Geometrization} but for the naive extended Hecke stacks, using the decomposition \eqref{eq:naivedecomp} to rigorously justify that all of the proofs go through verbatim.

In particular, the assignment $V \mapsto T^{e}_{V}$ gives a monoidal functor 
\begin{equation*}
    \tn{Rep}_{\Lambda}(\widehat{G} \rtimes Q) \to \tn{End}_{\tn{Rep}_{\Lambda}(Q)}(\mc{D}_{\tn{lis}}(\Bun_{G}^{e},\Lambda))^{B\Weil_{F}}.
\end{equation*}

Observe that base-changing $\Hck_{Z \subset G}^{e,\{0,1\},\tn{naive}}$ along our fixed divisor $\Spd(C) \xrightarrow{y} \Div^{1}$ yields a Cartesian diagram
\[
\begin{tikzcd}
\Hck_{Z \subset G}^{e,\{0,1\},\tn{naive}} \times_{\Div^{1}}  \Spd(C) \arrow{r} \arrow{d} & \Hck_{Z \subset G}^{e,1} \arrow{d} \\
\Hck_{Z \subset G}^{e,\loc,\{0,1\},\tn{naive}} \times_{\Div^{1}} \Spd(C) \arrow{r}  & \Hck_{Z \subset G}^{e,\loc,1}
\end{tikzcd}
\]
and that we have an embedding
\begin{equation*}
    \Hck_{Z \subset G}^{e,\{0,1\},\tn{naive}} \times_{\Div^{1}}  \Spd(C) \to \Hck_{Z \subset G}^{e,1}
\end{equation*}
which realizes $\Hck_{Z \subset G}^{e,\{0,1\},\tn{naive}} \times_{\Div^{1}} \Spd(C)$ as the substack
\begin{equation}\label{eq:naiveincl}
     \bigsqcup_{\lambda \in \Hom_{F}(u,Z)} \Hck_{Z \subset G}^{e,1,(\lambda,\lambda)} \hookrightarrow \bigsqcup_{(\lambda,\lambda') \in \Hom_{F}(u,Z)^{2}} \Hck_{Z \subset G}^{e,1,(\lambda,\lambda')} = \Hck_{Z \subset G}^{e,1},
\end{equation}
and that this also holds for the local analogues of both stacks.

Using these maps, we obtain:

\begin{thm}[Extended Hecke action]\label{thm:extended-hecke-action}
\begin{enumerate}
\item{There is a $\Weil_F^{\bullet}$-linear monoidal natural transformation 
\begin{equation*}
\bigoplus_{\Hom_{F}(u,Z)}\Rep_{\Lambda}(({^L}G)^{\bullet}) \xrightarrow{T} \End_{\Lambda}(\mc{D}_{\tn{lis}}(\Bun_{G}^{e},\Lambda))^{B\Weil_F^{\bullet}},
\end{equation*}
where the monoidal structure on the left-hand side is given by the usual tensor product on summands in the same graded component and is zero otherwise.
}
\item{There is a commutative diagram in $\Cat^{\otimes,\ex, \mathrm{st}}_{\Lambda}$  
\[\begin{tikzcd}
	{ \bigoplus_{(\lambda_{0}, \lambda_{1}) \in \Hom_{F}(u, Z)^{2}} \tn{Rep}_{\Lambda}(\widehat{G/Z})^{\lambda_{0}/\lambda_{1}} } & {\End_{\Lambda}(\mc{D}_{\tn{lis}}(\Bun_{G}^{e},\Lambda)) } \\
	{\bigoplus_{\Hom_{F}(u,Z)}\Rep_{\Lambda}({^L}G)} & {\End_{\Lambda}(\mc{D}_{\tn{lis}}(\Bun_{G}^{e},\Lambda))^{B\Weil_{F}},}
	\arrow["{T^e}", from=1-1, to=1-2]
	\arrow["\oblv", from=2-1, to=1-1]
	\arrow["T"', from=2-1, to=2-2]
	\arrow["\oblv"', from=2-2, to=1-2]
\end{tikzcd}\]}
\end{enumerate}
where the ``$\tn{oblv}$'' map first restricts to $\bigoplus_{\Hom_{F}(u,Z)}\Rep_{\Lambda}(\widehat{G})$ and then embeds via the map $\lambda \mapsto (\lambda,\lambda)$ on the gradings.
\end{thm}

\begin{rque}
    Part (1) of the above theorem implies that there is a $W_{F}^{\bullet}$-linear action of the monoidal category $\Rep_{\Lambda}(({^L}G)^{\bullet})$ on $\mc{D}_{\tn{lis}}(\Bun_{G}^{e},\Lambda) = \prod_{\lambda \in \Hom_{F}(u,Z)} \mc{D}_{\tn{lis}}(\Bun_{G}^{e,\lambda},\Lambda)$ which preserves each factor in the product.
\end{rque}

\begin{proof}
    The first part is Proposition \ref{prop:naiveSA}, which, to recall, just breaks up the naive Hecke stack and applies the theory of \cite{Geometrization} to each component. 
    
    For the second part, note that one has separate Hecke operators associated to $\Hck_{Z \subset G}^{e,\loc,\{0,1\},\tn{naive}} \times_{\Div^{1}} \Spd(C)$ and $\Hck_{Z \subset G}^{e,\loc,\{0,1\},\tn{naive}}$ (and their global analogues), which fit into a commutative square
    \[\begin{tikzcd}
	{\bigoplus_{\Hom_{F}(u,Z)} \Rep_{\Lambda}\widehat{G}} & {\End_{\Lambda}(\mc{D}_{\tn{lis}}(\Bun_{G}^{e},\Lambda)) } \\
	{\bigoplus_{\Hom_{F}(u,Z)} \Rep_{\Lambda}{^L}G} & {\End_{\Lambda}(\mc{D}_{\tn{lis}}(\Bun_{G}^{e},\Lambda))^{B\Weil_{F}}}.
	\arrow[from=1-1, to=1-2]
	\arrow["\oblv", from=2-1, to=1-1]
	\arrow["T"', from=2-1, to=2-2]
	\arrow["\oblv"', from=2-2, to=1-2]
\end{tikzcd}\]
   The claimed commutativity holds by \cite[Theorem 1.4.4]{Hansen24} (which itself summarizes \cite{Geometrization}), after applying the decomposition \eqref{eq:naivedecomp} to reduce to the non-extended case.

   We deduce the result from the commutativity of
    \[\begin{tikzcd}
	{\bigoplus_{(\lambda_{0},\lambda_{1}) \in \Hom_{F}(u,Z)^{2}}\Rep_{\Lambda}(\widehat{G/Z})^{\lambda_{0}/\lambda_{1}}} & {\End_{\Lambda}(\mc{D}_{\tn{lis}}(\Bun_{G}^{e},\Lambda)) } \\
	{\bigoplus_{\Hom_{F}(u,Z)}\Rep_{\Lambda}({\widehat{G})}} & {\End_{\Lambda}(\mc{D}_{\tn{lis}}(\Bun_{G}^{e},\Lambda))}
	\arrow["T^{e}", from=1-1, to=1-2]
	\arrow[from=2-1, to=1-1]
	\arrow[from=2-1, to=2-2]
	\arrow["\tn{id}"', from=2-2, to=1-2],
\end{tikzcd}\]
which is immediate from the inclusion \eqref{eq:naiveincl} and the above construction of the operators $T^{e}_{V^{e}}$.
\end{proof}

 By considering the monoidal embeddings
 \begin{equation*}
   \Rep_{\Lambda}({^L}G) \hookrightarrow \bigoplus_{\Hom_{F}(u,Z)}\Rep_{\Lambda}{^L}G, \hspace{1mm} \Rep_{\Lambda}(\widehat{G/Z) } \xrightarrow{\eqref{eq:G/Zpullback}} \bigoplus_{(\lambda_{0},\lambda_{1}) \in \Hom_{F}(u,Z)^{2}}\Rep_{\Lambda}(\widehat{G/Z})^{\lambda_{0}/\lambda_{1}},
 \end{equation*}
 one obtains:

 \begin{corol}\label{cor:maindiag}
 \begin{enumerate}
 \item{There is a $\Weil_F^{\bullet}$-linear monoidal natural transformation 
\begin{equation*}
\Rep_{\Lambda}(({^L}G)^{\bullet}) \to \End_{\Lambda}(\mc{D}_{\tn{lis}}(\Bun_{G}^{e},\Lambda))^{B\Weil_F^{\bullet}};
\end{equation*}
}
    \item{There is a commutative diagram in $\Cat^{\otimes,\ex, \mathrm{st}}_{\Lambda}$  
\[\begin{tikzcd}
	{\Rep_{\Lambda}(\widehat{G/Z})} & {\End_{\Lambda}(\mc{D}_{\tn{lis}}(\Bun_{G}^{e},\Lambda)) } \\
	{\Rep_{\Lambda}({^L}G}) & {\End_{\Lambda}(\mc{D}_{\tn{lis}}(\Bun_{G}^{e},\Lambda))^{B\Weil_{F}}.}
	\arrow["{T^e}", from=1-1, to=1-2]
	\arrow["\oblv", from=2-1, to=1-1]
	\arrow["T"', from=2-1, to=2-2]
	\arrow["\oblv"', from=2-2, to=1-2]
\end{tikzcd}\]}

\end{enumerate}
 \end{corol}

\begin{thm}\label{thm:diagHckact}
     There is a commutative diagram in $\Cat^{\otimes,\ex, \mathrm{st}}_{\Lambda}$  
\[\begin{tikzcd}
	{\Rep_{\Lambda}(\widehat{G}^{e})} & {\End_{\Lambda}(\mc{D}_{\tn{lis}}(\Bun_{G}^{e},\Lambda)) } \\
	{\Rep_{\Lambda}({^L}G}) & {\End_{\Lambda}(\mc{D}_{\tn{lis}}(\Bun_{G}^{e},\Lambda))^{B\Weil_{F}}.}
	\arrow["{T^e}", from=1-1, to=1-2]
	\arrow["\oblv", from=2-1, to=1-1]
	\arrow["T"', from=2-1, to=2-2]
	\arrow["\oblv"', from=2-2, to=1-2]
\end{tikzcd}\]
Moreover, for $V^{e} \in \Rep_{\Lambda}(\widehat{G}^{e})$, the action of $T^{e}_{V^{e}}$ on $\mc{D}_{\tn{lis}}(\Bun_{G}^{e},\Lambda)$ preserves all
limits and colimits, and the full subcategories of compact objects.
 \end{thm}

 \begin{proof}
     Fix $Z \subset Z' \subset Z_{G}$ both finite. The construction of the Hecke action for $Z$ takes place on the closed and open substack
     \begin{equation*}
\Hck_{Z \subset G}^{e,1} = \bigsqcup_{(\lambda_{0},\lambda_{1}) \in \Hom_{F}(u,Z)^{2}} \Hck_{Z' \subset G}^{e,1,(\lambda_{0},\lambda_{1})} \hookrightarrow  \bigsqcup_{(\lambda_{0},\lambda_{1}) \in \Hom_{F}(u,Z')} \Hck_{Z' \subset G}^{e,1,(\lambda_{0},\lambda_{1})} = \Hck_{Z' \subset G}^{e,1},
     \end{equation*}
     in a way which is compatible with the decomposition 
     \begin{equation*}
\Bun_{Z \subset G}^{e} = \bigsqcup_{\lambda \in \Hom_{F}(u,Z)} \Bun_{G}^{e,\lambda} \hookrightarrow \bigsqcup_{\lambda \in \Hom_{F}(u,Z')} \Bun_{G}^{e,\lambda} = \Bun_{Z' \subset G}^{e}.
     \end{equation*}
     The first part of the theorem then follows by Corollary \ref{cor:maindiag}, Lemma \ref{lem:changeinZconv}, and taking the direct limit over all $Z$.

     The proof of the second part comes from the identical argument as in the proof of \cite[Theorem IX.2.2]{Geometrization} (for each finite $Z$).
 \end{proof}

 \subsection{Dependence on choices}
Continue the notation of the previous subsections. There were two choices made above: 
\begin{enumerate}
\item{The choice of system of untilts $\{y^{(E)}\}_{E/F}$ lifting a given $y$ required for the identification $\tn{Gr}_{t}(C) \xrightarrow{\sim} X_{*}(t)$;}
\item{The choice of section $s$ of the surjection
\begin{equation}\label{eq:sectionsurj}
     \Hom_{F}(\tilde{t}, \mathbf{T}) \times_{\Hom_{F}(u,\mathbf{T})} \Hom_{F}(u,Z) \to \Hom_{F}(u,Z).
 \end{equation}
}
\end{enumerate}

\begin{rque}
    Although the pinning of $G$ was required to define each $\tn{Gr}_{Z \subset G}^{e,(\lambda_{0},\lambda_{1})}$, the Hecke action constructed above only depends on $\Hck_{Z \subset G}^{e,1,(\lambda_{0},\lambda_{1})}$, whose definition is independent of this choice.
\end{rque}

\subsubsection{The divisors}\label{sec:choiceoflifts}
The choice of a fixed $y \in \Div^{1}(C)$ already implicitly appears in \cite{Geometrization} by using the \'{e}tale fundamental group of $\Div^{1}$. For fixed $y$, any two choices of systems of untilts $\{y_{1}^{(E)}\}_{E/F}$ and $\{y_{2}^{(E)}\}_{E/F}$ are conjugate under some unique $\sigma \in \Gal_{F}$. The difference between the two resulting trivializations $\psi_{2} \circ \psi_{1}^{-1}$ is translation by a unique element $t_{\sigma} \in t(X_{C}^{\circ})$ from which one obtains $dt_{\sigma} \in Z^{1}(X_{C}^{\circ},u)$.

\begin{lem}\label{lem:sigmaaut}
    The class $[dt_{\sigma}] \in H_{\tn{v}}^{1}(X_{C}^{\circ},u)$ is nontrivial. In fact, it is even nontrivial in $\tn{Coker}[H_{\tn{fppf}}^{1}(\D_{C},u) \to H_{\tn{fppf}}^{1}(\D_{C}^{\circ},u)]$.
\end{lem}

\begin{proof}
Combining the Kummer sequence with Shapiro's Lemma gives an identification
\begin{equation*}
H_{\tn{v}}^{1}(X_{C}^{\circ},u) \xrightarrow{\sim} \varprojlim \Z[\Gal_{E/F}]_{0},
\end{equation*}
and under this map $[dt_{\sigma}]$ maps to $[e] -[\sigma]$.

For the second claim, we use purity for closed immersions (\cite[Theorem 7.1.2]{CS24}) to identify $\tn{Coker}[H_{\tn{fppf}}^{1}(\D_{C},u) \to H_{\tn{fppf}}^{1}(\D_{C}^{\circ},u)]$ with $H^{0}(C, u(-1)) = \varprojlim \widehat{\Z}[\Gal_{E/F}]$ and observe that $[dt_{\sigma}]$ again maps to $[e] - [\sigma]$. 
\end{proof}

In particular, since $H_{\tn{v}}^{1}(X_{C},u) = H_{\tn{fppf}}^{1}(F,u) = 0$ by Proposition \ref{prop:ucohom} (using \cite[Th\'{e}or\`{e}me 11.4]{Fargues22} for the first equality), the class $[dt_{\sigma}]$ does not lift to a class defined over $X_{C}$ and, locally, does not lift to a class defined over $\D_{C}$.

The cocycle $dt_{\sigma}$ defines an automorphism of the gerbe $[\mc{T}_{C}^{\circ}/\tilde{t}_{C}]$, denoted by $z_{\sigma}$. 

\begin{lem}
    For $\lambda \in \Hom_{F}(u,Z)$ and $S \in \Perfd_{C}$, the $Z_{S}$-torsor obtained using $\psi_{2}$, denoted by $\Q^{(2)}_{\lambda,S}$, is related to $\Q^{(1)}_{\lambda,S}$ (defined using $\psi_{1}$) by the formula 
    \begin{equation*}
\Q^{(2)}_{\lambda,S} = z_{\sigma,*}(\Q^{(1)}_{\lambda,S}).
    \end{equation*}
\end{lem}

Lemma \ref{lem:sigmaaut} (and the sentence following it) implies there is no automorphism of $[\mc{T}_{C}/\tilde{t}_{C}]$, or even of $[\mc{T}_{C}^{\loc}/\tilde{t}_{C}]$, lifting $z_{\sigma}$, and therefore no obvious way to use $z_{\sigma}$ to define an isomorphism of the (local or global) Hecke stacks associated to the two distinct $\{y_{i}^{(E)}\}$. It is therefore not clear to us what effect the choice of $\{y^{(E)}\}_{E/F}$ has on the resulting Hecke action (Theorem \ref{thm:extended-hecke-action}).

\begin{rque}
It is interesting to note the parallel between the necessity of choosing the lifts $\{y^{(E)}\}_{E/F}$ for $y \in \Div^{1}(C)$ in this paper and the necessity of choosing a section of the map $V_{\ov{F}} \to V_{F}$ (where $V$ denotes the set of places of a global field $F$) in order to define the global Kaletha gerbe in \cite{Kaletha18b}.
\end{rque}

\subsubsection{The section $s$}
The construction of the extended Hecke stacks (Definitions \ref{defi:elochck} and \ref{defi:eglobhck}) relied on a choice of section $s$ for the surjection \eqref{eq:sectionsurj}.
 
 Using Shapiro's lemma, we may identify $\Hom_{F}(\tilde{t},\mathbf{T})$ with $X_{*}(\mathbf{T})_{\mathbb{Q}}$ and $\Hom_{F}(u,\mathbf{T})$ with the group $X_{*}(\mathbf{T})_{\mathbb{Q}/\Z}$. It follows that picking a section of $\mathbb{Q} \to \mathbb{Q}/\mathbb{Z}$ determines a section of 
 \begin{equation*}
\Hom_{F}(\tilde{t},\mathbf{T}) \to \Hom_{F}(u,\mathbf{T}),
 \end{equation*}
 and therefore, by restricting to $\Hom_{F}(u,Z)$, the desired section. We can thus make a ``distinguished'' choice of section by taking the one corresponding to the section $\mathbb{Q}/\Z \to \mathbb{Q}$ valued in $[0,1)$. However, we will show shortly that the resulting Hecke action (Theorem \ref{thm:extended-hecke-action}) does not depend on the section $s$ at all. 
 
Any two sections $s_{1}$ and $s_{2}$ are related by the formula $s_{2} = s_{1}c$, where $\Hom_{F}(u,Z_{G}) \xrightarrow{c} \Hom_{F}(t,\mathbf{T})$ is a function (a ``$1$-cochain''). Denote by $\prescript{i}{}\Hck^{e,-}_{Z \subset G}$ the corresponding extended (local or global) Hecke stacks for $i=1,2$. There is an isomorphism
\begin{equation*}
    \prescript{1}{}\Hck^{e,-,n}_{Z \subset G} \xrightarrow{\eta_{c}} \prescript{2}{}\Hck^{e,-,n}_{Z \subset G}
\end{equation*}
sending (for $n=1$) $(\cP_{0}, \cP_{1}, f)$ to $(\cP_{0}, \cP_{1}, f \circ c(\lambda_{1}/\lambda_{0})_{*}(\psi)),$ where by $f \circ c(\lambda_{1}/\lambda_{0})_{*}(\psi)$ we mean the composition (using Lemma \ref{lem:curlyT})
\begin{equation*}
    \Q^{(2)}_{\lambda_{1}/\lambda_{0},S} \cdot \cP_{0} = (\Q^{(1)}_{\lambda_{1}/\lambda_{0},S} \cdot c(\lambda_{1}/\lambda_{0})_{*}(\mc{T}_{S})) \cdot \cP_{0} \xrightarrow{c(\lambda_{1}/\lambda_{0})_{*}(\psi)}\Q^{(1)}_{\lambda_{1}/\lambda_{0},S}  \cdot \cP_{0} \xrightarrow{f} \cP_{1}.
\end{equation*}

\begin{lem}
\begin{enumerate}
    \item{The map $\eta_{c}$ is compatible with convolution.}
    \item{The map $\eta_{c}$ commutes with the projection maps $\prescript{i}{}\Hck^{e,\loc,1}_{Z \subset G} \to \Hck_{G/Z}^{\loc,1}$ from Proposition \ref{prop:Hckquot}.}
    \end{enumerate}
\end{lem}

\begin{proof}
    The first claim is a straightforward calculation. Unpacking the constructions, the second claim reduces to proving the identity
    \begin{equation*}
        \varphi^{(1)}_{\lambda_{1}/\lambda_{0},\cP_{0}} = \ov{c(\lambda_{1}/\lambda_{0})_{*}(\psi)} \circ \varphi^{(2)}_{\lambda_{1}/\lambda_{0},\cP_{0}},
    \end{equation*}
    where $\varphi^{(i)}_{\lambda_{1}/\lambda_{0},\cP_{0}}$ denotes the canonical isomorphism $\ov{\cP_{0}} \xrightarrow{\sim} \ov{\Q_{\lambda_{1}/\lambda_{0},S}^{(i)} \cdot \cP_{0}}$ from \eqref{eq:varphi}. The desired identity then follows from the equality
    \begin{equation*}
        s_{1}(\lambda_{1}/\lambda_{0})_{*}(\psi^{-1}) = c(\lambda_{1}/\lambda_{0})_{*}(\psi) \circ s_{2}(\lambda_{1}/\lambda_{0})_{*}(\psi^{-1}) = c(\lambda_{1}/\lambda_{0})_{*}(\psi) \circ [s_{1}(\lambda_{1}/\lambda_{0})c(\lambda_{1}/\lambda_{0})]_{*}(\psi^{-1} )
    \end{equation*}
    of isomorphisms of torsors from $\underline{\mathbf{T}_{S}}$ to $s_{1}(\lambda_{1}/\lambda_{0})_{*}(\mc{T}_{S}^{\loc,\circ})$. 
\end{proof}

Putting the above lemma together with basic functoriality properties gives:

\begin{prop}
  All of the maps in the commutative diagram in part (2) of Theorem \ref{thm:extended-hecke-action} remain the same for any choice of section $s$ of \eqref{eq:sectionsurj}.
\end{prop}
 